\newpage
\section*{Appendix 3 --- Requirement verification, technology comparisons, and architecture details}
\addcontentsline{toc}{section}{Appendix 3 --- Requirement verification, technology comparisons, and architecture details}

Tables that support Sections 4.1.5, 4.3.1, 5.2, 5.3.1, and 5.3.4 and are kept out of the main text for length. Ratings in the comparison tables use the three symbols of Section 3.3.2: $\bullet$ meets the criterion, $\ominus$ meets it in part, $\circ$ does not.

User requirements.

User requirements --- explorer:

{\small\setstretch{1}\renewcommand{\arraystretch}{1.1}
\begin{longtable}{|>{\raggedright\arraybackslash}p{\dimexpr 0.114\textwidth-2\tabcolsep-\arrayrulewidth\relax}|>{\raggedright\arraybackslash}p{\dimexpr 0.744\textwidth-2\tabcolsep-\arrayrulewidth\relax}|>{\raggedright\arraybackslash}p{\dimexpr 0.137\textwidth-2\tabcolsep-\arrayrulewidth\relax}|}
\caption{User requirements of explorers.} \\
\hline
\textbf{ID} & \textbf{User requirement} & \textbf{Priority} \\ \hline
\endfirsthead
\hline
\textbf{ID} & \textbf{User requirement} & \textbf{Priority} \\ \hline
\endhead
UR-E01 & As an explorer, I want to create an account or sign in with Google, recover my account if I forget my password, and return to my own progress on every visit, so that my exploration is kept safe. & Must \\ \hline
UR-E02 & As an explorer, I want to browse the occupations available, read what each involves at every level, including its unfavourable aspects, and try the entry level of any of them without prerequisite, so that I can survey possibilities before committing to one. & Must \\ \hline
UR-E03 & As an explorer who does not know where to begin, I want optional suggestions of occupations to start with, and a description of myself that explains them, so that a blank choice does not stop me. & Should \\ \hline
UR-E04 & As an explorer, I want to be able to read my role, team, and working context before each scenario begins, so that my decisions are situated rather than abstract. & Must \\ \hline
UR-E05 & As an explorer, I want to perform tasks that practitioners would recognise and whose objective I can understand, so that what I learn reflects the occupation. & Must \\ \hline
UR-E06 & As an explorer, I want to respond in the forms the work itself takes --- choosing a course of action, ordering steps, prioritising tasks, or writing an answer such as an email or a code review --- so that I experience how the work is done, not only what it is about. & Must \\ \hline
UR-E07 & As an explorer, I want hints when I am stuck, at a small cost to my best possible score, so that I can continue while still learning. & Should \\ \hline
UR-E08 & As an explorer, I want my response judged as an expert in the occupation would judge it, probed further when my answer is unclear, and returned as a result, an explanation, and guidance, so that I can trust the judgment and know how to improve. & Must \\ \hline
UR-E09 & As an explorer, I want follow-up questions and feedback to arrive without long waits, and to see progress while I wait, so that the experience keeps its momentum. & Must \\ \hline
UR-E10 & As an explorer, I want to play on my phone, over a weak connection, and with assistive technology if I need it, without installing anything, so that I can explore wherever I am. & Must \\ \hline
UR-E11 & As an explorer, I want to resume where I stopped after an interruption or a system failure, and to be told promptly if something has failed, so that I never lose work I have done. & Must \\ \hline
UR-E12 & As an explorer, I want to know which content and judgments are produced by AI, so that I can weigh them accordingly. & Must \\ \hline
UR-E13 & As an explorer, I want to choose scenarios within a level, have the competencies I demonstrate recorded, and open higher levels as they grow, without replaying the same scenario being enough to inflate them, so that my progress reflects what I can do rather than time spent. & Must \\ \hline
UR-E14 & As an explorer, I want to start an adjacent occupation at a level my competencies support, so that I need not restart from the entry level when exploring related work. & Should \\ \hline
UR-E15 & As an explorer, I want to see the career path across the levels of an occupation --- the role held at each level --- so that I understand where the occupation could lead me. & Should \\ \hline
UR-E16 & As an explorer, I want to rate an occupation after completing it and see how others rated it, without anyone's individual rating being exposed, so that I can compare my impression with a broader one. & Should \\ \hline
UR-E17 & As an explorer, I want a private record of what I have attempted and demonstrated, so that my exploration accumulates into something I can reflect on. & Should \\ \hline
UR-E18 & As an explorer, I want my individual responses and results never to be shown to others or passed to employers, and my data to be deleted when I ask, so that trying an occupation carries no risk to me. & Must \\ \hline
UR-E19 & As an explorer, I want short side quests I can answer on the spot for bonus points, so that I can engage with an occupation in small steps. & Could \\ \hline
UR-E20 & As an explorer, I want to send feedback about the whole app, so that problems and ideas reach the team. & Should \\ \hline
\end{longtable}
}

User requirements --- domain expert:

{\small\setstretch{1}\renewcommand{\arraystretch}{1.1}
\begin{longtable}{|>{\raggedright\arraybackslash}p{\dimexpr 0.114\textwidth-2\tabcolsep-\arrayrulewidth\relax}|>{\raggedright\arraybackslash}p{\dimexpr 0.744\textwidth-2\tabcolsep-\arrayrulewidth\relax}|>{\raggedright\arraybackslash}p{\dimexpr 0.137\textwidth-2\tabcolsep-\arrayrulewidth\relax}|}
\caption{User requirements of domain experts.} \\
\hline
\textbf{ID} & \textbf{User requirement} & \textbf{Priority} \\ \hline
\endfirsthead
\hline
\textbf{ID} & \textbf{User requirement} & \textbf{Priority} \\ \hline
\endhead
UR-D01 & As a domain expert, I want to write a scenario's context with AI assistance and set the allowed task types, the score of each option, the grading rubric, hints, random events, and side quests directly in the system, with my authorship recorded, so that the content I own is traceable to me and needs no developer to enter. & Must \\ \hline
UR-D02 & As a domain expert, I want every published scenario to keep the required facts and option scores exactly as I set them, so that the system never misrepresents the work in my name. & Must \\ \hline
UR-D03 & As a domain expert, I want to review scenarios written by another expert of my occupation against shared criteria and approve them, request changes, or reject them, so that only content experts judge accurate, and AI output kept within the limits they set, reaches explorers. & Must \\ \hline
UR-D04 & As a domain expert, I want to see where the automated judge disagrees with expert judgments, so that I can refine the rubric. & Should \\ \hline
UR-D05 & As a domain expert, I want to revalidate published content when practice changes, so that what explorers see stays current. & Should \\ \hline
UR-D06 & As a domain expert, I want to create and keep up to date the profile of an occupation once the content administrator has created it --- levels, level descriptions, future roles, downsides, salary, generation rules, competencies, and adjacency --- so that progression and adjacency reflect how the work is organised in Vietnam. & Should \\ \hline
UR-D07 & As a domain expert, I want to rate explorers' responses and generated scenarios blind, independently of other experts, so that the automated judge and the generator can be measured against expert judgment. & Should \\ \hline
UR-D08 & As the author of a scenario, I want to play the draft game before it is published, without my play being counted as explorer data, so that I can confirm it presents my content correctly. & Must \\ \hline
UR-D09 & As the author of a scenario, I want to add sample answers with expert scores while preparing it, so that the automated judge can be checked against my judgment. & Should \\ \hline
\end{longtable}
}

User requirements --- content administrator:

{\small\setstretch{1}\renewcommand{\arraystretch}{1.1}
\begin{longtable}{|>{\raggedright\arraybackslash}p{\dimexpr 0.114\textwidth-2\tabcolsep-\arrayrulewidth\relax}|>{\raggedright\arraybackslash}p{\dimexpr 0.744\textwidth-2\tabcolsep-\arrayrulewidth\relax}|>{\raggedright\arraybackslash}p{\dimexpr 0.137\textwidth-2\tabcolsep-\arrayrulewidth\relax}|}
\caption{User requirements of content administrators.} \\
\hline
\textbf{ID} & \textbf{User requirement} & \textbf{Priority} \\ \hline
\endfirsthead
\hline
\textbf{ID} & \textbf{User requirement} & \textbf{Priority} \\ \hline
\endhead
UR-C01 & As a content administrator, I want to add a new occupation and track it through authoring, review, publication, and retirement without asking developers, so that catalogue growth is managed. & Must \\ \hline
UR-C02 & As a content administrator, I want to arrange how approved content is presented in gameplay --- choosing task types from those the domain expert allows, characters, and scenes --- without altering its substance, so that explorers receive it in a coherent form. & Must \\ \hline
UR-C03 & As a content administrator, I want to publish a scenario only after it has been approved and play-tested by its author, so that nothing unvalidated is released. & Must \\ \hline
UR-C04 & As a content administrator, I want to see coverage and play statistics for my occupations, so that I can identify outdated or weak content and arrange its revalidation. & Should \\ \hline
UR-C05 & As a content administrator, I want to replace an outdated scenario by removing the old one first and then opening a new one, so that explorers never meet superseded content. & Should \\ \hline
UR-C06 & As a content administrator, I want to open a scenario for each task and assign a domain expert other than its author to review it by a deadline, so that every scenario is reviewed independently and on time. & Must \\ \hline
UR-C07 & As a content administrator, I want to import occupational reference data from O*NET and ESCO, so that occupations start from an established source. & Should \\ \hline
UR-C08 & As a content administrator, I want to define the criteria domain experts use to review scenarios, so that every review follows the same standard. & Must \\ \hline
\end{longtable}
}

User requirements --- system administrator:

{\small\setstretch{1}\renewcommand{\arraystretch}{1.1}
\begin{longtable}{|>{\raggedright\arraybackslash}p{\dimexpr 0.114\textwidth-2\tabcolsep-\arrayrulewidth\relax}|>{\raggedright\arraybackslash}p{\dimexpr 0.744\textwidth-2\tabcolsep-\arrayrulewidth\relax}|>{\raggedright\arraybackslash}p{\dimexpr 0.137\textwidth-2\tabcolsep-\arrayrulewidth\relax}|}
\caption{User requirements of system administrators.} \\
\hline
\textbf{ID} & \textbf{User requirement} & \textbf{Priority} \\ \hline
\endfirsthead
\hline
\textbf{ID} & \textbf{User requirement} & \textbf{Priority} \\ \hline
\endhead
UR-S01 & As a system administrator, I want to configure which models generate and which assess, keep the two separate, and change provider without touching any authored content, so that assessment stays independent and the platform is not locked to one vendor. & Must \\ \hline
UR-S02 & As a system administrator, I want to monitor usage, cost, and failures, re-run failed AI jobs, and keep the cost of each session within budget, so that the service remains operable on grant funding. & Must \\ \hline
UR-S03 & As a system administrator, I want to manage every account, explorers included --- grant staff their roles and occupations, view activity logs, and disable or delete accounts --- so that each person can do exactly their part and misuse can be stopped. & Must \\ \hline
UR-S04 & As a system administrator, I want to review feedback sent about the app, so that recurring problems are found and fixed. & Should \\ \hline
\end{longtable}
}

Business model and organisation.

{\small\setstretch{1}\renewcommand{\arraystretch}{1.1}
\begin{longtable}{|>{\raggedright\arraybackslash}p{\dimexpr 0.171\textwidth-2\tabcolsep-\arrayrulewidth\relax}|>{\raggedright\arraybackslash}p{\dimexpr 0.159\textwidth-2\tabcolsep-\arrayrulewidth\relax}|>{\raggedright\arraybackslash}p{\dimexpr 0.301\textwidth-2\tabcolsep-\arrayrulewidth\relax}|>{\raggedright\arraybackslash}p{\dimexpr 0.193\textwidth-2\tabcolsep-\arrayrulewidth\relax}|>{\raggedright\arraybackslash}p{\dimexpr 0.171\textwidth-2\tabcolsep-\arrayrulewidth\relax}|}
\caption{Staged roadmap from university project to national public resource.} \\
\hline
\textbf{Stage} & \textbf{Indicative period} & \textbf{Activities} & \textbf{Resources and funding} & \textbf{Gate to next stage} \\ \hline
\endfirsthead
\hline
\textbf{Stage} & \textbf{Indicative period} & \textbf{Activities} & \textbf{Resources and funding} & \textbf{Gate to next stage} \\ \hline
\endhead
0. University project & Present phase, to the end of 2026 & Core system for the IT field; provisional seeds prepared by the team from authoritative public sources under advisor guidance; Stage A evaluation & Project team; advisor; host university; model costs self-funded or covered by cloud credits & Working system; Stage A targets met \\ \hline
1. Capstone and first pilot & Capstone phase, 2027 & Complete product; first expert panel for three pilot roles (back-end engineering, DevOps and cloud engineering, and business analysis); pilot with volunteer IT students; Stage B evaluation & University research or innovation support; cloud credits; competitions; donations & RP-1 and RP-2 thresholds met for the pilot roles \\ \hline
2. Incubation and multi-site pilot & 2027--2028 & Operating organisation established; expert panels extended through industry partners; pilots with partner schools and university career centres in Ho Chi Minh City & Edtech grants; company sponsorship; donations; in-kind expert consultation & Multi-site evidence of use and decision clarity; expert coverage of the IT roles \\ \hline
3. Public-\allowbreak{}sector partnership & 2028--2029 & Evidence presented to the city Department of Education and Training, then to the Ministry; alignment with the successor to Scheme 522; a second occupational field begins & Public programme funding; edtech grants; company sponsorship; donations & A formal cooperation agreement \\ \hline
4. National public resource & 2030 onward & Deployment through the public education system; further fields; annual expert revalidation & Public and philanthropic funding; company sponsorship; donations & Sustained operation \\ \hline
\end{longtable}
}

{\small\setstretch{1}\renewcommand{\arraystretch}{1.1}
\begin{longtable}{|>{\raggedright\arraybackslash}p{\dimexpr 0.182\textwidth-2\tabcolsep-\arrayrulewidth\relax}|>{\raggedright\arraybackslash}p{\dimexpr 0.273\textwidth-2\tabcolsep-\arrayrulewidth\relax}|>{\raggedright\arraybackslash}p{\dimexpr 0.336\textwidth-2\tabcolsep-\arrayrulewidth\relax}|>{\raggedright\arraybackslash}p{\dimexpr 0.204\textwidth-2\tabcolsep-\arrayrulewidth\relax}|}
\caption{Proposed organisational structure for sustained operation.} \\
\hline
\textbf{Unit} & \textbf{Positions (indicative number)} & \textbf{Responsibilities} & \textbf{Corresponding user group or constraint} \\ \hline
\endfirsthead
\hline
\textbf{Unit} & \textbf{Positions (indicative number)} & \textbf{Responsibilities} & \textbf{Corresponding user group or constraint} \\ \hline
\endhead
Product and Engineering & Engineering lead (1); full-stack engineers (2); AI engineer (1); DevOps and system administrator (1) & Platform development and operation; AI model integration, configuration, and cost control; infrastructure and security & System administrator; DC-01, DC-02, DC-05; NFR-11 \\ \hline
Content and Occupational Expertise & Content lead (1); assessment and rubric specialist (1); domain leads, one per occupational domain (part-\allowbreak{}time); domain experts (part-time, three to four per role) & Occupation catalogue; authoring of seeds, rubrics, and keyed options by domain experts, and their review of generated scenarios; presentation, review, and publication of occupations; annual revalidation & Domain expert; content administrator; BR-05, BR-07 \\ \hline
Research and Evaluation & Research lead (1); evaluation analyst (1) & Construction of the expert reference set; offline evaluation; subsequent user studies & Evaluation process; RP-1, RP-2 \\ \hline
Partnerships and Outreach & Partnerships manager (1); school and institution outreach coordinator (1) & Liaison with the Ministry, schools, and universities; grant applications and reporting; recruitment of domain experts & --- \\ \hline
Data Protection and Compliance & Data protection officer (1, part-time initially) & Consent; regulatory compliance & DC-03; policies in Section 4.1.4 \\ \hline
\end{longtable}
}

Comparison of technology choices.

{\small\setstretch{1}\renewcommand{\arraystretch}{1.1}
\begin{longtable}{|>{\raggedright\arraybackslash}p{\dimexpr 0.171\textwidth-2\tabcolsep-\arrayrulewidth\relax}|>{\raggedright\arraybackslash}p{\dimexpr 0.320\textwidth-2\tabcolsep-\arrayrulewidth\relax}|>{\raggedright\arraybackslash}p{\dimexpr 0.312\textwidth-2\tabcolsep-\arrayrulewidth\relax}|>{\raggedright\arraybackslash}p{\dimexpr 0.192\textwidth-2\tabcolsep-\arrayrulewidth\relax}|}
\caption{Comparison of frontend application frameworks.} \\
\hline
\textbf{Criterion} & \textbf{React 19 + Vite, single-page application (selected)} & \textbf{Next.js, server rendering\cite{r169}} & \textbf{Vue 3 + Vite, single-page application\cite{r170}} \\ \hline
\endfirsthead
\hline
\textbf{Criterion} & \textbf{React 19 + Vite, single-page application (selected)} & \textbf{Next.js, server rendering\cite{r169}} & \textbf{Vue 3 + Vite, single-page application\cite{r170}} \\ \hline
\endhead
Benefit the framework adds & $\circ$ no server rendering, which no screen needs (FR-01) & $\bullet$ server rendering and search indexing, unused behind sign-in & $\circ$ none beyond the selected option \\ \hline
TC3 Personal data & $\bullet$ static files hold no personal data and can sit on a content delivery network & $\ominus$ a rendering server handles personal data and must also be hosted in Vietnam & $\bullet$ static files \\ \hline
TC4 Weak devices and networks & $\bullet$ code split on demand; prototype entry about 195 KB & $\ominus$ less script, but each page waits on a server round trip & $\bullet$ small runtime \\ \hline
TC9 Feasibility & $\bullet$ prototype screens and 3D map already in React; largest ecosystem\cite{r134} & $\ominus$ same components, but server and client boundaries throughout the code & $\circ$ every screen rewritten; React Three Fiber replaced by TresJS\cite{r171} \\ \hline
\end{longtable}
}

{\small\setstretch{1}\renewcommand{\arraystretch}{1.1}
\begin{longtable}{|>{\raggedright\arraybackslash}p{\dimexpr 0.193\textwidth-2\tabcolsep-\arrayrulewidth\relax}|>{\raggedright\arraybackslash}p{\dimexpr 0.360\textwidth-2\tabcolsep-\arrayrulewidth\relax}|>{\raggedright\arraybackslash}p{\dimexpr 0.231\textwidth-2\tabcolsep-\arrayrulewidth\relax}|>{\raggedright\arraybackslash}p{\dimexpr 0.210\textwidth-2\tabcolsep-\arrayrulewidth\relax}|}
\caption{Comparison of occupation map renderings.} \\
\hline
\textbf{Criterion} & \textbf{3D view with an SVG view (selected)} & \textbf{3D view only} & \textbf{2D view only, in SVG} \\ \hline
\endfirsthead
\hline
\textbf{Criterion} & \textbf{3D view with an SVG view (selected)} & \textbf{3D view only} & \textbf{2D view only, in SVG} \\ \hline
\endhead
Progression metaphor (FR-05, FR-07) & $\bullet$ a galaxy of planets to travel between & $\bullet$ same & $\ominus$ a flat map, with a weaker sense of travel \\ \hline
TC4 Weak devices & $\bullet$ SVG view used automatically when WebGL is missing or memory is low & $\circ$ excludes devices without WebGL; GPU work every frame\cite{r139} & $\bullet$ lightest \\ \hline
TC5 Accessibility & $\bullet$ every planet a focusable, named element in the SVG view\cite{r140} & $\circ$ a canvas exposes nothing to a screen reader & $\bullet$ same as the selected option \\ \hline
TC9 Feasibility & $\ominus$ two views drawn from one data model & $\bullet$ exists in the prototype & $\ominus$ a new view; 3D work discarded \\ \hline
\end{longtable}
}

{\small\setstretch{1}\renewcommand{\arraystretch}{1.1}
\begin{longtable}{|>{\raggedright\arraybackslash}p{\dimexpr 0.171\textwidth-2\tabcolsep-\arrayrulewidth\relax}|>{\raggedright\arraybackslash}p{\dimexpr 0.291\textwidth-2\tabcolsep-\arrayrulewidth\relax}|>{\raggedright\arraybackslash}p{\dimexpr 0.186\textwidth-2\tabcolsep-\arrayrulewidth\relax}|>{\raggedright\arraybackslash}p{\dimexpr 0.347\textwidth-2\tabcolsep-\arrayrulewidth\relax}|}
\caption{Comparison of styling approaches.} \\
\hline
\textbf{Criterion} & \textbf{Tailwind CSS 4 (selected)} & \textbf{CSS Modules\cite{r172}} & \textbf{CSS-in-JS with a runtime, such as styled-\allowbreak{}components\cite{r173}} \\ \hline
\endfirsthead
\hline
\textbf{Criterion} & \textbf{Tailwind CSS 4 (selected)} & \textbf{CSS Modules\cite{r172}} & \textbf{CSS-in-JS with a runtime, such as styled-\allowbreak{}components\cite{r173}} \\ \hline
\endhead
TC4 Download & $\bullet$ only the classes used; no runtime\cite{r141} & $\bullet$ static CSS; no runtime & $\ominus$ styles generated in the browser by script \\ \hline
TC9 Feasibility & $\bullet$ used throughout the prototype & $\ominus$ styles rewritten & $\ominus$ styles rewritten \\ \hline
\end{longtable}
}

{\small\setstretch{1}\renewcommand{\arraystretch}{1.1}
\begin{longtable}{|>{\raggedright\arraybackslash}p{\dimexpr 0.175\textwidth-2\tabcolsep-\arrayrulewidth\relax}|>{\raggedright\arraybackslash}p{\dimexpr 0.286\textwidth-2\tabcolsep-\arrayrulewidth\relax}|>{\raggedright\arraybackslash}p{\dimexpr 0.200\textwidth-2\tabcolsep-\arrayrulewidth\relax}|>{\raggedright\arraybackslash}p{\dimexpr 0.334\textwidth-2\tabcolsep-\arrayrulewidth\relax}|}
\caption{Comparison of client state libraries.} \\
\hline
\textbf{Criterion} & \textbf{zustand (selected)} & \textbf{Redux Toolkit\cite{r174}} & \textbf{React context with useReducer} \\ \hline
\endfirsthead
\hline
\textbf{Criterion} & \textbf{zustand (selected)} & \textbf{Redux Toolkit\cite{r174}} & \textbf{React context with useReducer} \\ \hline
\endhead
TC2 Unsent turns kept & $\bullet$ persistence middleware\cite{r142} & $\ominus$ needs an added persistence library & $\circ$ written by hand \\ \hline
TC1 Re-\allowbreak{}rendering & $\bullet$ components subscribe to the part they use & $\bullet$ selectors & $\ominus$ every consumer re-renders on each change \\ \hline
TC4 Download & $\bullet$ very small & $\ominus$ larger & $\bullet$ nothing added \\ \hline
TC9 Feasibility & $\bullet$ used in the prototype & $\ominus$ migration & $\ominus$ persistence and subscriptions written by hand \\ \hline
\end{longtable}
}

{\small\setstretch{1}\renewcommand{\arraystretch}{1.1}
\begin{longtable}{|>{\raggedright\arraybackslash}p{\dimexpr 0.238\textwidth-2\tabcolsep-\arrayrulewidth\relax}|>{\raggedright\arraybackslash}p{\dimexpr 0.259\textwidth-2\tabcolsep-\arrayrulewidth\relax}|>{\raggedright\arraybackslash}p{\dimexpr 0.245\textwidth-2\tabcolsep-\arrayrulewidth\relax}|>{\raggedright\arraybackslash}p{\dimexpr 0.253\textwidth-2\tabcolsep-\arrayrulewidth\relax}|}
\caption{Comparison of routing approaches.} \\
\hline
\textbf{Criterion} & \textbf{React Router 7 (selected)} & \textbf{TanStack Router\cite{r175}} & \textbf{Screens switched by application state, without a router} \\ \hline
\endfirsthead
\hline
\textbf{Criterion} & \textbf{React Router 7 (selected)} & \textbf{TanStack Router\cite{r175}} & \textbf{Screens switched by application state, without a router} \\ \hline
\endhead
TC4 Code loaded on demand & $\bullet$ lazy routes & $\bullet$ lazy routes & $\ominus$ split by hand \\ \hline
Areas separated by role (FR-41) & $\bullet$ nested layouts & $\bullet$ nested layouts & $\ominus$ checks written by hand \\ \hline
TC2 Run resumed on any device & $\bullet$ route loaders fetch the run's state\cite{r143} & $\bullet$ route loaders & $\circ$ no address for a run \\ \hline
TC9 Feasibility & $\bullet$ used in the prototype & $\ominus$ migration, in exchange for type-\allowbreak{}checked routes & $\circ$ \\ \hline
\end{longtable}
}

{\small\setstretch{1}\renewcommand{\arraystretch}{1.1}
\begin{longtable}{|>{\raggedright\arraybackslash}p{\dimexpr 0.264\textwidth-2\tabcolsep-\arrayrulewidth\relax}|>{\raggedright\arraybackslash}p{\dimexpr 0.235\textwidth-2\tabcolsep-\arrayrulewidth\relax}|>{\raggedright\arraybackslash}p{\dimexpr 0.256\textwidth-2\tabcolsep-\arrayrulewidth\relax}|>{\raggedright\arraybackslash}p{\dimexpr 0.240\textwidth-2\tabcolsep-\arrayrulewidth\relax}|}
\caption{Comparison of result delivery mechanisms.} \\
\hline
\textbf{Criterion} & \textbf{Server-sent events (selected)} & \textbf{WebSocket\cite{r176}} & \textbf{Polling} \\ \hline
\endfirsthead
\hline
\textbf{Criterion} & \textbf{Server-sent events (selected)} & \textbf{WebSocket\cite{r176}} & \textbf{Polling} \\ \hline
\endhead
TC1 Result shown when ready & $\bullet$ pushed & $\bullet$ pushed & $\ominus$ delayed by the polling interval \\ \hline
TC2, TC4 Dropped connections & $\bullet$ the browser reconnects and resumes from the last event\cite{r144} & $\ominus$ reconnection and replay written by hand & $\bullet$ each request independent \\ \hline
Proxies and load balancers & $\bullet$ plain HTTP & $\ominus$ needs support for upgrading the connection & $\bullet$ plain HTTP \\ \hline
TC8 Load at 5,000 concurrent explorers & $\bullet$ one open stream each & $\bullet$ one connection each & $\circ$ a request every few seconds from each explorer \\ \hline
Direction needed: server to client only & $\bullet$ exactly that & $\ominus$ a two-way channel, half unused & $\ominus$ client asks repeatedly \\ \hline
\end{longtable}
}

{\small\setstretch{1}\renewcommand{\arraystretch}{1.1}
\begin{longtable}{|>{\raggedright\arraybackslash}p{\dimexpr 0.210\textwidth-2\tabcolsep-\arrayrulewidth\relax}|>{\raggedright\arraybackslash}p{\dimexpr 0.302\textwidth-2\tabcolsep-\arrayrulewidth\relax}|>{\raggedright\arraybackslash}p{\dimexpr 0.263\textwidth-2\tabcolsep-\arrayrulewidth\relax}|>{\raggedright\arraybackslash}p{\dimexpr 0.220\textwidth-2\tabcolsep-\arrayrulewidth\relax}|}
\caption{Comparison of staff interface foundations.} \\
\hline
\textbf{Criterion} & \textbf{React Aria, unstyled primitives (selected)\cite{r146}} & \textbf{MUI component library\cite{r177}} & \textbf{Components built by hand} \\ \hline
\endfirsthead
\hline
\textbf{Criterion} & \textbf{React Aria, unstyled primitives (selected)\cite{r146}} & \textbf{MUI component library\cite{r177}} & \textbf{Components built by hand} \\ \hline
\endhead
TC5 Accessible dialogs, comboboxes, and grids & $\bullet$ accessible behaviour built in & $\bullet$ accessible components & $\circ$ where hand-built interfaces most often fail \\ \hline
Fit with the Tailwind styling & $\bullet$ unstyled, styled with Tailwind & $\ominus$ its own styling system beside Tailwind & $\bullet$ \\ \hline
TC9 Feasibility & $\ominus$ a new library with a small surface & $\ominus$ a new library with a large surface & $\circ$ the most code to write and test \\ \hline
\end{longtable}
}

{\small\setstretch{1}\renewcommand{\arraystretch}{1.1}
\begin{longtable}{|>{\raggedright\arraybackslash}p{\dimexpr 0.237\textwidth-2\tabcolsep-\arrayrulewidth\relax}|>{\raggedright\arraybackslash}p{\dimexpr 0.370\textwidth-2\tabcolsep-\arrayrulewidth\relax}|>{\raggedright\arraybackslash}p{\dimexpr 0.203\textwidth-2\tabcolsep-\arrayrulewidth\relax}|>{\raggedright\arraybackslash}p{\dimexpr 0.186\textwidth-2\tabcolsep-\arrayrulewidth\relax}|}
\caption{Comparison of localisation approaches.} \\
\hline
\textbf{Criterion} & \textbf{Typed message table with the Intl API (selected)} & \textbf{i18next\cite{r178}} & \textbf{Strings written in components} \\ \hline
\endfirsthead
\hline
\textbf{Criterion} & \textbf{Typed message table with the Intl API (selected)} & \textbf{i18next\cite{r178}} & \textbf{Strings written in components} \\ \hline
\endhead
TC5 Missing strings caught & $\bullet$ an unknown key fails compilation; a test checks completeness & $\ominus$ type checking needs extra setup & $\circ$ not detectable \\ \hline
English added later as data & $\bullet$ & $\bullet$ & $\circ$ every screen edited \\ \hline
TC4 Download & $\bullet$ no library & $\ominus$ a library added & $\bullet$ \\ \hline
TC9 Feasibility & $\bullet$ used in the prototype & $\ominus$ migration & $\circ$ \\ \hline
\end{longtable}
}

{\small\setstretch{1}\renewcommand{\arraystretch}{1.1}
\begin{longtable}{|>{\raggedright\arraybackslash}p{\dimexpr 0.205\textwidth-2\tabcolsep-\arrayrulewidth\relax}|>{\raggedright\arraybackslash}p{\dimexpr 0.294\textwidth-2\tabcolsep-\arrayrulewidth\relax}|>{\raggedright\arraybackslash}p{\dimexpr 0.297\textwidth-2\tabcolsep-\arrayrulewidth\relax}|>{\raggedright\arraybackslash}p{\dimexpr 0.199\textwidth-2\tabcolsep-\arrayrulewidth\relax}|}
\caption{Comparison of backend frameworks.} \\
\hline
\textbf{Criterion} & \textbf{NestJS 11 on Node.js (selected)} & \textbf{Spring Boot, Java\cite{r179}} & \textbf{FastAPI, Python\cite{r181}} \\ \hline
\endfirsthead
\hline
\textbf{Criterion} & \textbf{NestJS 11 on Node.js (selected)} & \textbf{Spring Boot, Java\cite{r179}} & \textbf{FastAPI, Python\cite{r181}} \\ \hline
\endhead
TC1 Waiting on models and open streams & $\bullet$ event loop; a waiting request costs memory, not a thread & $\ominus$ a thread per request by default; virtual threads\cite{r180} or a reactive stack for thousands of streams & $\bullet$ asynchronous by design \\ \hline
TC6 Same rules as the browser & $\bullet$ imports game-core & $\circ$ rules rewritten in Java; two copies to keep equal & $\circ$ rules rewritten in Python \\ \hline
Modules, role guards, queue and stream integration & $\bullet$ first-party\cite{r149} & $\bullet$ mature equivalents & $\ominus$ assembled from separate libraries \\ \hline
Access to AI libraries & $\ominus$ provider SDKs in TypeScript; statistics offline in Python & $\ominus$ fewer AI libraries & $\bullet$ the Python ecosystem \\ \hline
TC9 Feasibility & $\bullet$ the prototype API is NestJS; the team writes TypeScript & $\ominus$ a new language and framework & $\ominus$ a new language for the API \\ \hline
\end{longtable}
}

{\small\setstretch{1}\renewcommand{\arraystretch}{1.1}
\begin{longtable}{|>{\raggedright\arraybackslash}p{\dimexpr 0.268\textwidth-2\tabcolsep-\arrayrulewidth\relax}|>{\raggedright\arraybackslash}p{\dimexpr 0.242\textwidth-2\tabcolsep-\arrayrulewidth\relax}|>{\raggedright\arraybackslash}p{\dimexpr 0.196\textwidth-2\tabcolsep-\arrayrulewidth\relax}|>{\raggedright\arraybackslash}p{\dimexpr 0.289\textwidth-2\tabcolsep-\arrayrulewidth\relax}|}
\caption{Comparison of databases.} \\
\hline
\textbf{Criterion} & \textbf{PostgreSQL with Prisma (selected)} & \textbf{MongoDB\cite{r182}} & \textbf{MySQL 8\cite{r183}} \\ \hline
\endfirsthead
\hline
\textbf{Criterion} & \textbf{PostgreSQL with Prisma (selected)} & \textbf{MongoDB\cite{r182}} & \textbf{MySQL 8\cite{r183}} \\ \hline
\endhead
Rules across records enforced by the database (FR-25, FR-33, FR-40) & $\bullet$ foreign keys, checks, partial unique indexes\cite{r151} & $\ominus$ no foreign keys; rules kept in application code & $\ominus$ foreign keys and checks; no partial indexes, functional indexes as a workaround \\ \hline
Publication and its log in one transaction (NFR-16) & $\bullet$ transactions & $\bullet$ multi-\allowbreak{}document transactions\cite{r182} & $\bullet$ transactions \\ \hline
Nested scenario content (NFR-10, NFR-15) & $\bullet$ indexable jsonb\cite{r152} & $\bullet$ native documents & $\ominus$ JSON type, indexed through functional indexes \\ \hline
Statistics above a threshold (NFR-06) & $\bullet$ window functions and materialised views & $\ominus$ aggregation pipeline & $\ominus$ window functions; no materialised views \\ \hline
TC9 Feasibility & $\bullet$ the prototype schema, 19 tables & $\circ$ data re-modelled & $\ominus$ supported by Prisma; migrations rewritten \\ \hline
\end{longtable}
}

{\small\setstretch{1}\renewcommand{\arraystretch}{1.1}
\begin{longtable}{|>{\raggedright\arraybackslash}p{\dimexpr 0.278\textwidth-2\tabcolsep-\arrayrulewidth\relax}|>{\raggedright\arraybackslash}p{\dimexpr 0.248\textwidth-2\tabcolsep-\arrayrulewidth\relax}|>{\raggedright\arraybackslash}p{\dimexpr 0.219\textwidth-2\tabcolsep-\arrayrulewidth\relax}|>{\raggedright\arraybackslash}p{\dimexpr 0.250\textwidth-2\tabcolsep-\arrayrulewidth\relax}|}
\caption{Comparison of job queues.} \\
\hline
\textbf{Criterion} & \textbf{Redis with BullMQ (selected)} & \textbf{RabbitMQ\cite{r184}} & \textbf{pg-boss on PostgreSQL\cite{r185}} \\ \hline
\endfirsthead
\hline
\textbf{Criterion} & \textbf{Redis with BullMQ (selected)} & \textbf{RabbitMQ\cite{r184}} & \textbf{pg-boss on PostgreSQL\cite{r185}} \\ \hline
\endhead
Retries with backoff, priorities, failed jobs listed (NFR-04, FR-43) & $\bullet$ built in\cite{r154} & $\ominus$ priorities and dead-letter queues built in; delayed retries configured by hand & $\bullet$ retries with backoff, priority and dead-letter queues \\ \hline
Rate limits per provider (NFR-11) & $\bullet$ per-queue limiter & $\ominus$ written by hand & $\bullet$ queue policies \\ \hline
Results to every API replica (NFR-02) & $\bullet$ publish--subscribe & $\bullet$ exchanges & $\ominus$ publish--subscribe through the database \\ \hline
Atomic cost counter per session (NFR-11) & $\bullet$ atomic counters & $\circ$ needs a separate store & $\bullet$ transactional update \\ \hline
Load on the system of record at Stage 4 & $\bullet$ none; Redis holds nothing that cannot be rebuilt & $\bullet$ none & $\ominus$ about one hundred jobs per second written to PostgreSQL \\ \hline
TC8 Services to operate & $\ominus$ one more service & $\ominus$ one more service, plus a store for counters & $\bullet$ none added \\ \hline
TC9 Feasibility & $\bullet$ first-party NestJS integration\cite{r149}; declared in the prototype & $\ominus$ new & $\ominus$ new \\ \hline
\end{longtable}
}

{\small\setstretch{1}\renewcommand{\arraystretch}{1.1}
\begin{longtable}{|>{\raggedright\arraybackslash}p{\dimexpr 0.171\textwidth-2\tabcolsep-\arrayrulewidth\relax}|>{\raggedright\arraybackslash}p{\dimexpr 0.339\textwidth-2\tabcolsep-\arrayrulewidth\relax}|>{\raggedright\arraybackslash}p{\dimexpr 0.281\textwidth-2\tabcolsep-\arrayrulewidth\relax}|>{\raggedright\arraybackslash}p{\dimexpr 0.204\textwidth-2\tabcolsep-\arrayrulewidth\relax}|}
\caption{Comparison of authentication solutions.} \\
\hline
\textbf{Criterion} & \textbf{The platform's own module (selected)} & \textbf{Keycloak, self-hosted\cite{r186}} & \textbf{Managed service, such as Firebase Authentication\cite{r187}} \\ \hline
\endfirsthead
\hline
\textbf{Criterion} & \textbf{The platform's own module (selected)} & \textbf{Keycloak, self-hosted\cite{r186}} & \textbf{Managed service, such as Firebase Authentication\cite{r187}} \\ \hline
\endhead
TC3 Identity data in Vietnam & $\bullet$ in the platform's database in Vietnam & $\bullet$ hosted in Vietnam & $\circ$ held by the provider abroad \\ \hline
Roles scoped to occupations (FR-41) & $\bullet$ in the application & $\ominus$ roles in Keycloak; occupation scope still checked in the application & $\ominus$ same \\ \hline
Security baseline (NFR-17) & $\ominus$ to be completed to ASVS level 2\cite{r157}: reset, idle expiry, second factor for staff & $\bullet$ reset, sessions, second factor built in & $\bullet$ built in \\ \hline
TC8 Operation & $\bullet$ no extra service & $\circ$ a separate identity server to run, upgrade, and secure & $\ominus$ no server; one vendor for every account \\ \hline
TC9 Feasibility & $\bullet$ registration, Google sign-in, rotating refresh tokens exist in the prototype & $\ominus$ migration & $\ominus$ migration \\ \hline
\end{longtable}
}

{\small\setstretch{1}\renewcommand{\arraystretch}{1.1}
\begin{longtable}{|>{\raggedright\arraybackslash}p{\dimexpr 0.258\textwidth-2\tabcolsep-\arrayrulewidth\relax}|>{\raggedright\arraybackslash}p{\dimexpr 0.363\textwidth-2\tabcolsep-\arrayrulewidth\relax}|>{\raggedright\arraybackslash}p{\dimexpr 0.214\textwidth-2\tabcolsep-\arrayrulewidth\relax}|>{\raggedright\arraybackslash}p{\dimexpr 0.159\textwidth-2\tabcolsep-\arrayrulewidth\relax}|}
\caption{Comparison of validation libraries.} \\
\hline
\textbf{Criterion} & \textbf{Zod 4 (selected)} & \textbf{class-\allowbreak{}validator\cite{r188}} & \textbf{JSON Schema with Ajv\cite{r189}} \\ \hline
\endfirsthead
\hline
\textbf{Criterion} & \textbf{Zod 4 (selected)} & \textbf{class-\allowbreak{}validator\cite{r188}} & \textbf{JSON Schema with Ajv\cite{r189}} \\ \hline
\endhead
TC6 One schema for browser, API, and worker & $\bullet$ plain TypeScript that runs anywhere & $\ominus$ decorated classes, suited to the server & $\bullet$ language-\allowbreak{}neutral \\ \hline
Schemas for model output (FR-11) & $\bullet$ exported as JSON Schema for providers' structured output\cite{r190} & $\ominus$ needs another library to export & $\bullet$ native \\ \hline
Types in the code & $\bullet$ inferred from the schema & $\ominus$ classes double as types & $\ominus$ types generated separately \\ \hline
Read by the offline Python harness & $\bullet$ through the JSON Schema export & $\circ$ & $\bullet$ \\ \hline
TC9 Feasibility & $\bullet$ the prototype validates game data with Zod & $\ominus$ & $\ominus$ \\ \hline
\end{longtable}
}

{\small\setstretch{1}\renewcommand{\arraystretch}{1.1}
\begin{longtable}{|>{\raggedright\arraybackslash}p{\dimexpr 0.294\textwidth-2\tabcolsep-\arrayrulewidth\relax}|>{\raggedright\arraybackslash}p{\dimexpr 0.253\textwidth-2\tabcolsep-\arrayrulewidth\relax}|>{\raggedright\arraybackslash}p{\dimexpr 0.281\textwidth-2\tabcolsep-\arrayrulewidth\relax}|>{\raggedright\arraybackslash}p{\dimexpr 0.167\textwidth-2\tabcolsep-\arrayrulewidth\relax}|}
\caption{Comparison of code organisation.} \\
\hline
\textbf{Criterion} & \textbf{Nx monorepo (selected)} & \textbf{Separate repositories, shared package published} & \textbf{pnpm workspaces alone} \\ \hline
\endfirsthead
\hline
\textbf{Criterion} & \textbf{Nx monorepo (selected)} & \textbf{Separate repositories, shared package published} & \textbf{pnpm workspaces alone} \\ \hline
\endhead
TC6 One copy of game-core and the schemas & $\bullet$ & $\ominus$ versions published and kept in step & $\bullet$ \\ \hline
TC4 Staff code kept out of the explorer's download & $\bullet$ module-\allowbreak{}boundary rules enforced\cite{r158} & $\bullet$ separated by construction & $\ominus$ not enforced \\ \hline
Only what changed is built and tested & $\bullet$ affected tasks and caching\cite{r158} & $\ominus$ per repository & $\ominus$ by hand \\ \hline
TC9 Feasibility & $\bullet$ the prototype is an Nx workspace & $\circ$ repository split & $\ominus$ Nx already runs on pnpm workspaces \\ \hline
\end{longtable}
}

{\small\setstretch{1}\renewcommand{\arraystretch}{1.1}
\begin{longtable}{|>{\raggedright\arraybackslash}p{\dimexpr 0.153\textwidth-2\tabcolsep-\arrayrulewidth\relax}|>{\raggedright\arraybackslash}p{\dimexpr 0.300\textwidth-2\tabcolsep-\arrayrulewidth\relax}|>{\raggedright\arraybackslash}p{\dimexpr 0.186\textwidth-2\tabcolsep-\arrayrulewidth\relax}|>{\raggedright\arraybackslash}p{\dimexpr 0.356\textwidth-2\tabcolsep-\arrayrulewidth\relax}|}
\caption{Comparison of hosting options.} \\
\hline
\textbf{Criterion} & \textbf{Containers hosted in Vietnam (selected)} & \textbf{Cloud region abroad, such as Singapore} & \textbf{Serverless functions} \\ \hline
\endfirsthead
\hline
\textbf{Criterion} & \textbf{Containers hosted in Vietnam (selected)} & \textbf{Cloud region abroad, such as Singapore} & \textbf{Serverless functions} \\ \hline
\endhead
TC3 Data in Vietnam & $\bullet$ system of record in Vietnam; only de-\allowbreak{}identified text leaves & $\circ$ every account becomes a cross-\allowbreak{}border transfer\cite{r99} & $\ominus$ only where the platform has a region in Vietnam \\ \hline
TC1 Request path & $\bullet$ long-\allowbreak{}running processes & $\bullet$ long-\allowbreak{}running processes & $\ominus$ cold starts on the request path (Section 5.3.1) \\ \hline
Open streams and queue workers & $\bullet$ & $\bullet$ & $\circ$ short-lived functions; streams and workers need other services \\ \hline
TC8 Operation & $\ominus$ the team runs containers; a university server at Stage 1 & $\ominus$ same & $\bullet$ scaling managed by the platform \\ \hline
Moving host & $\bullet$ portable services only & $\bullet$ & $\circ$ tied to one platform's interfaces \\ \hline
\end{longtable}
}

{\small\setstretch{1}\renewcommand{\arraystretch}{1.1}
\begin{longtable}{|>{\raggedright\arraybackslash}p{\dimexpr 0.190\textwidth-2\tabcolsep-\arrayrulewidth\relax}|>{\raggedright\arraybackslash}p{\dimexpr 0.328\textwidth-2\tabcolsep-\arrayrulewidth\relax}|>{\raggedright\arraybackslash}p{\dimexpr 0.191\textwidth-2\tabcolsep-\arrayrulewidth\relax}|>{\raggedright\arraybackslash}p{\dimexpr 0.287\textwidth-2\tabcolsep-\arrayrulewidth\relax}|}
\caption{Comparison of ways to hold generation to expert content.} \\
\hline
\textbf{Criterion} & \textbf{Whole seed in the prompt, structured output, rule checks (selected)} & \textbf{Retrieval over the expert material} & \textbf{Fine-tuning a model on expert scenarios} \\ \hline
\endfirsthead
\hline
\textbf{Criterion} & \textbf{Whole seed in the prompt, structured output, rule checks (selected)} & \textbf{Retrieval over the expert material} & \textbf{Fine-tuning a model on expert scenarios} \\ \hline
\endhead
Fidelity to expert content (RP-1, NFR-15) & $\bullet$ the whole seed in context; canon shown verbatim; names and numbers checked\cite{r120,r121} & $\ominus$ retrieval can miss a required fact & $\circ$ facts held in weights; required facts not guaranteed \\ \hline
Data needed in this phase & $\bullet$ one approved seed & $\bullet$ the same material, indexed & $\circ$ many expert-\allowbreak{}written scenarios, which do not yet exist (Section 4.2.1) \\ \hline
Seed revised by an expert (BR-14) & $\bullet$ takes effect at once & $\bullet$ re-indexed & $\circ$ model retrained \\ \hline
TC7 Provider replaceable & $\bullet$ any model with a long enough context & $\bullet$ & $\circ$ tied to the tuned model \\ \hline
TC8 Cost per call & $\ominus$ long prompts, reduced by prompt caching & $\ominus$ shorter prompts, plus an index to run & $\ominus$ cheaper calls, plus training \\ \hline
\end{longtable}
}

{\small\setstretch{1}\renewcommand{\arraystretch}{1.1}
\begin{longtable}{|>{\raggedright\arraybackslash}p{\dimexpr 0.219\textwidth-2\tabcolsep-\arrayrulewidth\relax}|>{\raggedright\arraybackslash}p{\dimexpr 0.372\textwidth-2\tabcolsep-\arrayrulewidth\relax}|>{\raggedright\arraybackslash}p{\dimexpr 0.222\textwidth-2\tabcolsep-\arrayrulewidth\relax}|>{\raggedright\arraybackslash}p{\dimexpr 0.182\textwidth-2\tabcolsep-\arrayrulewidth\relax}|}
\caption{Comparison of ways to check generated text at run time.} \\
\hline
\textbf{Criterion} & \textbf{Deterministic rules at run time, log review by another family afterwards (selected)} & \textbf{A guard model on every turn\cite{r199}} & \textbf{Prompt instructions only} \\ \hline
\endfirsthead
\hline
\textbf{Criterion} & \textbf{Deterministic rules at run time, log review by another family afterwards (selected)} & \textbf{A guard model on every turn\cite{r199}} & \textbf{Prompt instructions only} \\ \hline
\endhead
NFR-01 Time added while waiting & $\bullet$ milliseconds & $\circ$ a second model call on every turn & $\bullet$ none \\ \hline
NFR-15 Invented names and numbers & $\bullet$ every name and number must occur in the seed or the explorer's text & $\ominus$ judged by a model, with its own errors & $\circ$ not checked \\ \hline
Follow-up limits and leaked rubric text & $\bullet$ counted and matched by the engine & $\ominus$ & $\circ$ \\ \hline
Tone, plausibility, subtle drift & $\ominus$ caught afterwards by the log review & $\bullet$ caught before display & $\circ$ \\ \hline
TC6 One rule set & $\bullet$ the same checks in \textit{game-\allowbreak{}core}, client and server & $\circ$ a separate model to configure & $\circ$ \\ \hline
\end{longtable}
}

{\small\setstretch{1}\renewcommand{\arraystretch}{1.1}
\begin{longtable}{|>{\raggedright\arraybackslash}p{\dimexpr 0.189\textwidth-2\tabcolsep-\arrayrulewidth\relax}|>{\raggedright\arraybackslash}p{\dimexpr 0.331\textwidth-2\tabcolsep-\arrayrulewidth\relax}|>{\raggedright\arraybackslash}p{\dimexpr 0.213\textwidth-2\tabcolsep-\arrayrulewidth\relax}|>{\raggedright\arraybackslash}p{\dimexpr 0.261\textwidth-2\tabcolsep-\arrayrulewidth\relax}|}
\caption{Comparison of scoring designs.} \\
\hline
\textbf{Criterion} & \textbf{Panel of three judges from three families, median (selected)} & \textbf{One strong judge} & \textbf{Fine-tuned scoring model} \\ \hline
\endfirsthead
\hline
\textbf{Criterion} & \textbf{Panel of three judges from three families, median (selected)} & \textbf{One strong judge} & \textbf{Fine-tuned scoring model} \\ \hline
\endhead
Agreement with experts & $\ominus$ reaches the best single judge's level without choosing it in advance\cite{r123}; to be measured & $\ominus$ depends on choosing the right model & $\bullet$ highest on essay benchmarks\cite{r125}, after training on large human-\allowbreak{}scored sets \\ \hline
Data needed in this phase & $\bullet$ rubric only & $\bullet$ rubric only & $\circ$ large expert-\allowbreak{}scored sets, unavailable before Stage B \\ \hline
Self-\allowbreak{}preference and single-\allowbreak{}model bias (DC-01) & $\bullet$ three families, none the generator's\cite{r128} & $\ominus$ one family; its biases unchecked & $\ominus$ biases of its training data \\ \hline
Signal for expert review (AI-3) & $\bullet$ disagreement between judges flags a case & $\circ$ confidence separates right from wrong only weakly\cite{r117} & $\circ$ same \\ \hline
TC2 One provider down & $\bullet$ the other tier of the same family & $\ominus$ switching family changes the judge & $\circ$ the self-hosted model must stay up \\ \hline
TC8, NFR-02 Cost and time & $\ominus$ three calls in parallel; time of the slowest & $\bullet$ one call & $\bullet$ cheap calls after training \\ \hline
\end{longtable}
}

{\small\setstretch{1}\renewcommand{\arraystretch}{1.1}
\begin{longtable}{|>{\raggedright\arraybackslash}p{\dimexpr 0.165\textwidth-2\tabcolsep-\arrayrulewidth\relax}|>{\raggedright\arraybackslash}p{\dimexpr 0.460\textwidth-2\tabcolsep-\arrayrulewidth\relax}|>{\raggedright\arraybackslash}p{\dimexpr 0.211\textwidth-2\tabcolsep-\arrayrulewidth\relax}|>{\raggedright\arraybackslash}p{\dimexpr 0.159\textwidth-2\tabcolsep-\arrayrulewidth\relax}|}
\caption{Comparison of ways to obtain well-formed model output.} \\
\hline
\textbf{Criterion} & \textbf{Schema-\allowbreak{}constrained output where the provider offers it, validation, one retry (selected)} & \textbf{Prompted format, then repair by a second model} & \textbf{Tool (function) calling} \\ \hline
\endfirsthead
\hline
\textbf{Criterion} & \textbf{Schema-\allowbreak{}constrained output where the provider offers it, validation, one retry (selected)} & \textbf{Prompted format, then repair by a second model} & \textbf{Tool (function) calling} \\ \hline
\endhead
Output always parses (DC-02) & $\bullet$ guaranteed by the provider, or checked and retried & $\ominus$ repaired after the fact & $\ominus$ arguments may still fail the schema \\ \hline
Judgment unchanged by the fix & $\bullet$ a failed output is discarded, never edited & $\circ$ the repairing model can alter the level & $\bullet$ \\ \hline
Every candidate model & $\ominus$ strict schemas on the judges' providers; Qwen 3.8 Flash outside its Singapore region\cite{r200}; JSON only on DeepSeek\cite{r201} & $\bullet$ & $\ominus$ \\ \hline
NFR-01, NFR-02 Time & $\bullet$ no extra call when valid & $\circ$ a second call on failure & $\bullet$ \\ \hline
\end{longtable}
}

{\small\setstretch{1}\renewcommand{\arraystretch}{1.1}
\begin{longtable}{|>{\raggedright\arraybackslash}p{\dimexpr 0.193\textwidth-2\tabcolsep-\arrayrulewidth\relax}|>{\raggedright\arraybackslash}p{\dimexpr 0.290\textwidth-2\tabcolsep-\arrayrulewidth\relax}|>{\raggedright\arraybackslash}p{\dimexpr 0.159\textwidth-2\tabcolsep-\arrayrulewidth\relax}|>{\raggedright\arraybackslash}p{\dimexpr 0.171\textwidth-2\tabcolsep-\arrayrulewidth\relax}|>{\raggedright\arraybackslash}p{\dimexpr 0.182\textwidth-2\tabcolsep-\arrayrulewidth\relax}|}
\caption{Comparison of ways to reach model providers.} \\
\hline
\textbf{Criterion} & \textbf{Own adapter layer, OpenAI-\allowbreak{}compatible format (selected)} & \textbf{Each provider's SDK in domain code} & \textbf{LiteLLM proxy\cite{r191}} & \textbf{OpenRouter, hosted router\cite{r192}} \\ \hline
\endfirsthead
\hline
\textbf{Criterion} & \textbf{Own adapter layer, OpenAI-\allowbreak{}compatible format (selected)} & \textbf{Each provider's SDK in domain code} & \textbf{LiteLLM proxy\cite{r191}} & \textbf{OpenRouter, hosted router\cite{r192}} \\ \hline
\endhead
TC7 Provider replaced by configuration (DC-02, NFR-12) & $\bullet$ & $\circ$ code changed per provider & $\bullet$ & $\bullet$ \\ \hline
TC7 Judge families kept apart (DC-01) & $\bullet$ registry in the gateway & $\ominus$ checks scattered & $\ominus$ still written in the application & $\ominus$ still written in the application \\ \hline
TC3 Path of the data & $\bullet$ direct to each provider & $\bullet$ direct & $\bullet$ through a self-hosted proxy & $\circ$ through an added intermediary abroad \\ \hline
Structured output from every provider & $\ominus$ native adapters where compatibility is limited: Anthropic's layer is not meant for production and ignores strict schemas\cite{r195}; Google's ignores unsupported parameters\cite{r196} & $\bullet$ full native features & $\ominus$ depends on the proxy's mapping & $\ominus$ depends on the router's mapping \\ \hline
Local model as a fallback (idea only, Section 5.3.6) & $\bullet$ Ollama and vLLM serve the same API\cite{r193,r194} & $\ominus$ another SDK & $\bullet$ & $\circ$ cannot reach a local model \\ \hline
TC8 Operation & $\bullet$ a library inside the API and worker & $\bullet$ & $\ominus$ an extra Python service to run & $\ominus$ no service, but a paid intermediary \\ \hline
\end{longtable}
}

{\small\setstretch{1}\renewcommand{\arraystretch}{1.1}
\begin{longtable}{|>{\raggedright\arraybackslash}p{\dimexpr 0.350\textwidth-2\tabcolsep-\arrayrulewidth\relax}|>{\raggedright\arraybackslash}p{\dimexpr 0.391\textwidth-2\tabcolsep-\arrayrulewidth\relax}|>{\raggedright\arraybackslash}p{\dimexpr 0.254\textwidth-2\tabcolsep-\arrayrulewidth\relax}|}
\caption{Comparison of narrator candidates.} \\
\hline
\textbf{Criterion} & \textbf{qwen3.\allowbreak{}8-\allowbreak{}flash (provisional narrator)} & \textbf{deepseek-\allowbreak{}flash} \\ \hline
\endfirsthead
\hline
\textbf{Criterion} & \textbf{qwen3.\allowbreak{}8-\allowbreak{}flash (provisional narrator)} & \textbf{deepseek-\allowbreak{}flash} \\ \hline
\endhead
NFR-08 Vietnamese & $\bullet$ 74.9, the higher of the two & $\ominus$ 72.1 \\ \hline
NFR-11 Price & $\bullet$ 0.15 / 0.47 & $\ominus$ 0.30 / 1.20 \\ \hline
Well-formed output & $\bullet$ strict schema, outside the Singapore region & $\ominus$ JSON mode; empty replies retried \\ \hline
DC-01 Family apart from every judge & $\bullet$ & $\bullet$ \\ \hline
Knowledge of IT work; time from Vietnam & measured in Section 7.3 & measured in Section 7.3 \\ \hline
\end{longtable}
}

{\small\setstretch{1}\renewcommand{\arraystretch}{1.1}
\begin{longtable}{|>{\raggedright\arraybackslash}p{\dimexpr 0.309\textwidth-2\tabcolsep-\arrayrulewidth\relax}|>{\raggedright\arraybackslash}p{\dimexpr 0.283\textwidth-2\tabcolsep-\arrayrulewidth\relax}|>{\raggedright\arraybackslash}p{\dimexpr 0.255\textwidth-2\tabcolsep-\arrayrulewidth\relax}|>{\raggedright\arraybackslash}p{\dimexpr 0.148\textwidth-2\tabcolsep-\arrayrulewidth\relax}|}
\caption{Comparison of offline evaluation harnesses.} \\
\hline
\textbf{Criterion} & \textbf{Inspect AI (selected)\cite{r129}} & \textbf{promptfoo\cite{r198}} & \textbf{Scripts written by the project} \\ \hline
\endfirsthead
\hline
\textbf{Criterion} & \textbf{Inspect AI (selected)\cite{r129}} & \textbf{promptfoo\cite{r198}} & \textbf{Scripts written by the project} \\ \hline
\endhead
Repeated runs, several graders, median & $\bullet$ built in & $\ominus$ repeats and model-\allowbreak{}graded rubrics; aggregation by hand & $\circ$ written by hand \\ \hline
Ordinal agreement statistics & $\ominus$ Krippendorff's alpha built in; quadratic weighted kappa added as a metric & $\circ$ computed outside the tool & $\ominus$ written with standard libraries \\ \hline
Custom scorer for the project's rubric & $\bullet$ scorer interface & $\bullet$ custom assertions & $\bullet$ \\ \hline
Logs for reviewing each judgment & $\bullet$ log viewer & $\bullet$ web viewer & $\ominus$ \\ \hline
Shares prompts and rubrics with the TypeScript product & $\ominus$ Python, reading the shared data files & $\bullet$ Node.js with YAML configuration & $\bullet$ \\ \hline
\end{longtable}
}

{\small\setstretch{1}\renewcommand{\arraystretch}{1.1}
\begin{longtable}{|>{\raggedright\arraybackslash}p{\dimexpr 0.171\textwidth-2\tabcolsep-\arrayrulewidth\relax}|>{\raggedright\arraybackslash}p{\dimexpr 0.313\textwidth-2\tabcolsep-\arrayrulewidth\relax}|>{\raggedright\arraybackslash}p{\dimexpr 0.184\textwidth-2\tabcolsep-\arrayrulewidth\relax}|>{\raggedright\arraybackslash}p{\dimexpr 0.186\textwidth-2\tabcolsep-\arrayrulewidth\relax}|>{\raggedright\arraybackslash}p{\dimexpr 0.141\textwidth-2\tabcolsep-\arrayrulewidth\relax}|}
\caption{Comparison of ways to use language models on explorer text.} \\
\hline
\textbf{Criterion} & \textbf{De-\allowbreak{}identify, then send to providers abroad, treated as a transfer (selected)} & \textbf{Send answers abroad as they are, under consent} & \textbf{Open-weight models served in Vietnam} & \textbf{No language model on explorer text} \\ \hline
\endfirsthead
\hline
\textbf{Criterion} & \textbf{De-\allowbreak{}identify, then send to providers abroad, treated as a transfer (selected)} & \textbf{Send answers abroad as they are, under consent} & \textbf{Open-weight models served in Vietnam} & \textbf{No language model on explorer text} \\ \hline
\endhead
DC-03, NFR-05 Personal data in Vietnam & $\ominus$ identifiers removed before sending; residual risk in free text & $\circ$ every answer leaves with its author's identity & $\bullet$ nothing leaves & $\bullet$ nothing leaves \\ \hline
RP-2, NFR-14 Scoring quality & $\bullet$ the strongest Vietnamese-\allowbreak{}capable models (Section 5.2.3) & $\bullet$ & $\ominus$ smaller models; quality to be shown & $\circ$ keyword scoring fails (Section 7.3.3) \\ \hline
TC8, TC9 Cost and feasibility now & $\bullet$ pay per call; no hardware & $\bullet$ & $\circ$ GPU hosting beyond the project's means & $\bullet$ \\ \hline
Legal obligations & $\ominus$ transfer impact assessment and notice & $\circ$ the same, for identified personal data, at scale & $\bullet$ none for transfer & $\bullet$ none for transfer \\ \hline
\end{longtable}
}

Architecture.

{\small\setstretch{1}\renewcommand{\arraystretch}{1.1}
\begin{longtable}{|>{\raggedright\arraybackslash}p{\dimexpr 0.310\textwidth-2\tabcolsep-\arrayrulewidth\relax}|>{\raggedright\arraybackslash}p{\dimexpr 0.214\textwidth-2\tabcolsep-\arrayrulewidth\relax}|>{\raggedright\arraybackslash}p{\dimexpr 0.470\textwidth-2\tabcolsep-\arrayrulewidth\relax}|}
\caption{Realisation of use cases by architectural components.} \\
\hline
\textbf{Use cases} & \textbf{Components} & \textbf{Mechanism} \\ \hline
\endfirsthead
\hline
\textbf{Use cases} & \textbf{Components} & \textbf{Mechanism} \\ \hline
\endhead
Register account; Manage account & Explorer application; Identity and access; Notification & Own authentication module; single-use reset links by email \\ \hline
Complete Get-\allowbreak{}to-\allowbreak{}know-\allowbreak{}me quiz; Build user portrait & Explorer application; Exploration & RIASEC matching in \textit{game-core} without AI; fixed portrait text for the highest-\allowbreak{}scoring type loaded from the database and labelled as self-\allowbreak{}reported (FR-04, DR-12) \\ \hline
Explore occupations; View occupation profile; View level description; View career path; Enter adjacent occupation & Explorer application; Exploration; Analytics; Game core & Map model from the published release; 3D or accessible 2D rendering; adjacency offers computed by \textit{game-core} from DR-02 and DR-03; ratings with suppression \\ \hline
Choose scenario; Complete side quest & Exploration; Game core & Enrolment and credit in the database; thresholds and caps evaluated by \textit{game-core} on the server; adjacency precomputed from imported data \\ \hline
Perform scenario; View scenario context; Get hint; Resolve activity; Answer follow-up question; Flag wrong content & Explorer application; Exploration; Assessment; AI orchestration; Model gateway & The assessed turn above; flags routed to the occupation's domain experts \\ \hline
Rate occupation; View personal record; Send app feedback & Exploration; Analytics; Operations & Ratings only after completion, aggregated on write; records rebuilt from events; feedback stored with app version \\ \hline
Manage occupation profile; Author scenario; Generate scenario content; Build reference set & Staff application; Content lifecycle; AI orchestration & Drafts with authorship; generation jobs with history and automatic validation; reference answers validated on entry \\ \hline
Review scenario; Apply review criteria; Revalidate scenario; Review flagged content; Play-test draft game & Staff application; Content lifecycle; Exploration; Notification; Scheduler & State machine with reviewer-\allowbreak{}is-\allowbreak{}not-\allowbreak{}author constraint; criteria versioned per review; play-test runs excluded from aggregates \\ \hline
Check grading reliability; Submit blind ratings & Staff application; Analytics; Assessment & Agreement by stratum from recorded judgments; blind ratings hidden from other experts and from AI scores \\ \hline
Manage occupation; Import reference data; Define review criteria; Open task scenario; Assign domain expert; Reassign overdue review & Staff application; Content lifecycle; Reference-\allowbreak{}data import; Identity and access; Scheduler & One active scenario per task enforced by the database; reviewer and deadline set on assignment, author refused; reminders and overdue notices by the scheduler \\ \hline
Design scenario presentation; Generate missing asset; Publish scenario; Remove scenario & Staff application; Content lifecycle; AI orchestration; Object store & Content text read-only during presentation; task types limited to the domain expert's allowed set; assets checked before entering the library; publication creates an immutable release; deletion retires and keeps history \\ \hline
Monitor occupation statistics; View coverage report; Request revalidation & Staff application; Analytics; Content lifecycle & Aggregates above threshold; revalidation sends a published scenario back to review while it stays playable \\ \hline
Manage user accounts; Assign occupation staff; Disable account; Delete account; View activity log & Staff application; Identity and access & Occupation-\allowbreak{}scoped roles; status checked on refresh and on every staff request; deletion keeps staff-\allowbreak{}authored records \\ \hline
Configure AI models; Monitor operation cost; Re-run failed job; Review app feedback & Staff application; Operations; Model gateway & Registry rejects a provider or family shared by the narrator and a judge or by two judges; cost ledger and alerts; failed-job re-run \\ \hline
\end{longtable}
}

{\small\setstretch{1}\renewcommand{\arraystretch}{1.1}
\begin{longtable}{|>{\raggedright\arraybackslash}p{\dimexpr 0.114\textwidth-2\tabcolsep-\arrayrulewidth\relax}|>{\raggedright\arraybackslash}p{\dimexpr 0.528\textwidth-2\tabcolsep-\arrayrulewidth\relax}|>{\raggedright\arraybackslash}p{\dimexpr 0.204\textwidth-2\tabcolsep-\arrayrulewidth\relax}|>{\raggedright\arraybackslash}p{\dimexpr 0.148\textwidth-2\tabcolsep-\arrayrulewidth\relax}|}
\caption{Realisation of non-\allowbreak{}functional requirements and design constraints by architectural mechanisms.} \\
\hline
\textbf{ID} & \textbf{Mechanism} & \textbf{Components} & \textbf{In prototype} \\ \hline
\endfirsthead
\hline
\textbf{ID} & \textbf{Mechanism} & \textbf{Components} & \textbf{In prototype} \\ \hline
\endhead
NFR-01 & Answer stored before any model call; narration and follow-up generated by a worker and pushed through the event stream; progress at once, status line after two seconds; interactive jobs before batch generation; bounded output length & Explorer application; Exploration; AI orchestration; Model gateway & No; follow-ups are pre-written \\ \hline
NFR-02 & Assessment job queued as soon as each answer is stored; judges in parallel, so latency is the slowest judge's, not the sum; per-call timeouts; workers scaled on queue depth & AI orchestration; Model gateway; workers & No \\ \hline
NFR-03 & Each turn stored as an event with an idempotent sequence number; local queue of unsent turns; resumption by folding stored events & Explorer application; Exploration; Game core & Partly; runs pinned to a content version, but actions sent only at the end \\ \hline
NFR-04 & Transactional outbox; retries with backoff, circuit breaker, failover within the role; notice within ten seconds; reconciliation & Exploration; AI orchestration; Model gateway; Scheduler & No \\ \hline
NFR-05 & Data inventory per data requirement; encryption in transit and at rest, including backups; de-\allowbreak{}identification before any model call; opaque job identifiers; providers set to no training and no retention; deletion by cascade & Identity and access; AI orchestration; Operations; database & Partly; deletion by cascade \\ \hline
NFR-06 & Aggregates only from Analytics, above a population threshold; play-tests excluded; no endpoint returns another person's record & Analytics & No \\ \hline
NFR-07 & Every model output carries provenance --- model, family, prompt version --- from which the interface renders an AI label & AI orchestration; Explorer application & No \\ \hline
NFR-08 & Typed interface strings with a completeness test; content fields carry a locale, so a language is added as data; Vietnamese only in this phase & Both applications; Content lifecycle & Partly; interface strings typed \\ \hline
NFR-09 & Validator refuses a scenario whose context lacks an explicit objective & Content lifecycle; Game core & No \\ \hline
NFR-10 & Content as data in releases, created through the staff application; \textit{game-core} interprets generic formats; map positions computed from data & Content lifecycle; Game core; Explorer application & Partly; content in the database, written by script \\ \hline
NFR-11 & Per-session cost ledger: estimate before, actual after each call; ceiling from settings; threshold alerts; a session that reaches the ceiling makes no further model calls and pauses as on provider failure & AI orchestration; Operations; queue store & No \\ \hline
NFR-12 & Model port with provider adapters; provider-\allowbreak{}neutral prompts and output schemas; a provider change re-runs the Stage A benchmark as a regression suite & Model gateway; AI orchestration & No \\ \hline
NFR-13 & Explorer bundle under a performance budget; on-demand 3D chunk and 2D fallback; turn payloads of a few kilobytes; event stream with polling fallback; releases cached by version at the edge; on-device timing & Explorer application; content delivery network & Partly; 3D map loaded on demand \\ \hline
NFR-14 & Every judgment records its judges; a change of any judge model recomputes agreement by stratum against the reference set & Analytics; AI orchestration; Operations & No \\ \hline
NFR-15 & Publication copies the approved version into a hashed release, compared with what is served; canon content served verbatim; run-time narration shown only after the rule checks, with the expert's default follow-up when a follow-up fails them; feedback notes redrafted until they pass; a model of another family reviews the logs & Content lifecycle; AI orchestration; Game core & Partly; releases immutable \\ \hline
NFR-16 & Lifecycle module as sole writer; transition and log in one transaction; a trigger refuses updates to published releases & Content lifecycle; database & No \\ \hline
NFR-17 & Argon2id; HTTPS with strict transport security; short-lived bearer access tokens held in memory, rotating refresh tokens, idle expiry; answers passed to models as delimited data, outputs schema-\allowbreak{}validated; prompt-\allowbreak{}injection tests in the regression suite, against the leading OWASP risk for language-\allowbreak{}model applications\cite{r166} & Identity and access; AI orchestration; reverse proxy & Partly; hashing, bearer tokens, rotation \\ \hline
NFR-18 & Stateless API replicas with health checks and rolling deployment; backups with point-\allowbreak{}in-\allowbreak{}time recovery; uptime monitoring; announced maintenance & All tiers & No \\ \hline
NFR-19 & Document-\allowbreak{}based interfaces; accessible 2D map; accessible primitives in the staff application; automated and manual audits & Both applications & No \\ \hline
DC-01 & Registry records each model's provider and family; configuration rejects either shared by the narrator and a judge, or by two judges; panel rechecked at job time & Model gateway; Operations; AI orchestration & No \\ \hline
DC-02 & Domain code depends only on the model port; providers are adapters selected per role & Model gateway & No \\ \hline
DC-03 & System of record in Vietnam; one egress path, through de-\allowbreak{}identification & Hosting; AI orchestration; Identity and access; Operations & No \\ \hline
DC-05 & Browser applications, no installation; static delivery; tolerant of short disconnections & Both applications & Yes \\ \hline
\end{longtable}
}

{\footnotesize\setstretch{1}\renewcommand{\arraystretch}{1.1}
\begin{longtable}{|>{\raggedright\arraybackslash}p{\dimexpr 0.158\textwidth-2\tabcolsep-\arrayrulewidth\relax}|>{\raggedright\arraybackslash}p{\dimexpr 0.148\textwidth-2\tabcolsep-\arrayrulewidth\relax}|>{\raggedright\arraybackslash}p{\dimexpr 0.148\textwidth-2\tabcolsep-\arrayrulewidth\relax}|>{\raggedright\arraybackslash}p{\dimexpr 0.128\textwidth-2\tabcolsep-\arrayrulewidth\relax}|>{\raggedright\arraybackslash}p{\dimexpr 0.108\textwidth-2\tabcolsep-\arrayrulewidth\relax}|>{\raggedright\arraybackslash}p{\dimexpr 0.305\textwidth-2\tabcolsep-\arrayrulewidth\relax}|}
\caption{Capacity estimate and topology by roadmap stage.} \\
\hline
\textbf{Stage} & \textbf{Assumed users} & \textbf{Peak concurrent explorers} & \textbf{API requests per second} & \textbf{Model calls per second} & \textbf{Topology} \\ \hline
\endfirsthead
\hline
\textbf{Stage} & \textbf{Assumed users} & \textbf{Peak concurrent explorers} & \textbf{API requests per second} & \textbf{Model calls per second} & \textbf{Topology} \\ \hline
\endhead
0. University project & Team and synthetic data & Under 10 & Under 5 & Under 1 & One host; all containers; hosted development database acceptable \\ \hline
1. Capstone and first pilot & 100 to 300 volunteers & About 50 & 5 to 10 & About 1 & University host in Vietnam; one API and one worker container; PostgreSQL with point-\allowbreak{}in-\allowbreak{}time recovery; Redis with append-only persistence; object store; uptime monitor; CDN for static files \\ \hline
2. Multi-site pilot & About 10,000 in partner schools & About 500, in class-sized bursts & 50 to 100 & About 8 & Cloud in Vietnam; two or more API replicas behind a load balancer; workers scaled on queue length; managed PostgreSQL with standby and read replica; managed Redis; staging \\ \hline
3 and 4. Public-\allowbreak{}sector partnership and national resource & Hundreds of thousands per year & About 5,000 & 500 to 1,000 & 80 to 90 & Two zones in Vietnam with autoscaling, PostgreSQL high availability, and a recovery region (deployment diagram) \\ \hline
\end{longtable}
}

{\small\setstretch{1}\renewcommand{\arraystretch}{1.1}
\begin{longtable}{|>{\raggedright\arraybackslash}p{\dimexpr 0.159\textwidth-2\tabcolsep-\arrayrulewidth\relax}|>{\raggedright\arraybackslash}p{\dimexpr 0.238\textwidth-2\tabcolsep-\arrayrulewidth\relax}|>{\raggedright\arraybackslash}p{\dimexpr 0.290\textwidth-2\tabcolsep-\arrayrulewidth\relax}|>{\raggedright\arraybackslash}p{\dimexpr 0.307\textwidth-2\tabcolsep-\arrayrulewidth\relax}|}
\caption{Scaling mechanism of each tier at Stage 4.} \\
\hline
\textbf{Tier} & \textbf{Load at Stage 4} & \textbf{How it scales} & \textbf{Why it can} \\ \hline
\endfirsthead
\hline
\textbf{Tier} & \textbf{Load at Stage 4} & \textbf{How it scales} & \textbf{Why it can} \\ \hline
\endhead
Content delivery network & All static files and published releases & Cache by version with long lifetimes & Releases are immutable, so caches never need purging and content reads never reach the database \\ \hline
API & 500 to 1,000 requests per second; about 5,000 open result streams & Replicas behind the load balancer, scaled on CPU and open connections; at least two per zone & Stateless: sessions travel in bearer tokens and results arrive through Redis publish--subscribe, so any replica can serve any explorer \\ \hline
Workers & 80 to 90 model calls per second; 250 to 700 in flight & Scaled on queue length; each worker runs many I/O-bound jobs concurrently; in-flight calls capped per provider & Jobs are independent and idempotent; the queue absorbs bursts\cite{r168} \\ \hline
Model providers & About 5,100 requests per minute; 7 to 8 million tokens per minute at about 1,500 tokens per call & Several providers per role within DC-01, so quotas add up; priorities: live grading and narration first, feedback notes next, authoring work last & The binding constraint is isolated in its own tier and limited without slowing the rest \\ \hline
PostgreSQL & Under 500 writes per second; about 0.3 TB of run data per season at about 60 KB per run & Vertical growth behind a connection pooler; read replica for analytics; run events partitioned by month & All transactions stay within one primary; heavy reads move to the replica; old partitions are compressed or archived \\ \hline
Redis & About 100 jobs per second, plus publish--subscribe messages & One primary and one replica & Holds nothing that cannot be rebuilt from PostgreSQL \\ \hline
\end{longtable}
}

Correspondence of sequence diagrams with methods.

Correspondence with the sequence diagrams. Every message of the sequence diagrams SD-01 to SD-12 is matched to a method, except the few listed under the open points below. The table gives the main entry and the chain behind it for each diagram.

{\small\setstretch{1}\renewcommand{\arraystretch}{1.1}
\begin{longtable}{|>{\raggedright\arraybackslash}p{\dimexpr 0.171\textwidth-2\tabcolsep-\arrayrulewidth\relax}|>{\raggedright\arraybackslash}p{\dimexpr 0.300\textwidth-2\tabcolsep-\arrayrulewidth\relax}|>{\raggedright\arraybackslash}p{\dimexpr 0.524\textwidth-2\tabcolsep-\arrayrulewidth\relax}|}
\caption{Entry method and chain of methods for each sequence diagram.} \\
\hline
\textbf{Diagram} & \textbf{Entry} & \textbf{Chain of methods} \\ \hline
\endfirsthead
\hline
\textbf{Diagram} & \textbf{Entry} & \textbf{Chain of methods} \\ \hline
\endhead
SD-01 Answer an activity & \texttt{ShiftRunController.\allowbreak{}takeTurn} & \texttt{RuleEngine.\allowbreak{}validateAnswer} $\rightarrow$ \texttt{AssessmentService.\allowbreak{}requestGrading} $\rightarrow$ \texttt{DeIdentifierService.\allowbreak{}strip} $\rightarrow$ \texttt{JudgePanelService.\allowbreak{}grade} and \texttt{combine} $\rightarrow$ \texttt{NarrationProcessor.\allowbreak{}process} $\rightarrow$ \texttt{OutputCheckService.\allowbreak{}verifyQuotes}, \texttt{checkNarration} $\rightarrow$ \texttt{CostLedgerService.\allowbreak{}record} $\rightarrow$ \texttt{StreamService.\allowbreak{}publish}; the feedback note by \texttt{FeedbackNoteProcessor.\allowbreak{}process} \\ \hline
SD-02 Finish a scenario & \texttt{ShiftRunController.\allowbreak{}finishRun} & \texttt{SummaryService.\allowbreak{}finish} $\rightarrow$ \texttt{ScoringService.\allowbreak{}groupPoints} $\rightarrow$ \texttt{PlayerLevelService.\allowbreak{}completeIfDone} $\rightarrow$ \texttt{SideQuestService.\allowbreak{}eligible} $\rightarrow$ \texttt{ProfileService.\allowbreak{}applyRun} \\ \hline
SD-03 Get-\allowbreak{}to-\allowbreak{}know-\allowbreak{}me quiz & \texttt{GetToKnowMeController.\allowbreak{}submitGetToKnowMe} & \texttt{GetToKnowMeService.\allowbreak{}submit} (self-\allowbreak{}reported RIASEC) $\rightarrow$ \texttt{SuggestionService.\allowbreak{}forPlayer} \\ \hline
SD-04 Access the platform & \texttt{AuthController.\allowbreak{}signIn}, \texttt{signInWithGoogle}, \texttt{register} & \texttt{PasswordService.\allowbreak{}verify}, \texttt{GoogleOAuthService.\allowbreak{}exchangeCode} $\rightarrow$ \texttt{AccountStore} $\rightarrow$ \texttt{TokenService.\allowbreak{}issue} $\rightarrow$ \texttt{AuditLogService.\allowbreak{}log} \\ \hline
SD-05 Generate scenario content & \texttt{GenerationController.\allowbreak{}requestGeneration} & \texttt{GenerationProcessor.\allowbreak{}process} $\rightarrow$ \texttt{ModelGatewayService.\allowbreak{}complete}; \texttt{SeedValidator.\allowbreak{}check} before \texttt{SeedController.\allowbreak{}submitScenario} \\ \hline
SD-06 Review a scenario & \texttt{ReviewController.\allowbreak{}decideReview} & \texttt{AccessPolicyService.\allowbreak{}isReviewerOf} $\rightarrow$ \texttt{ReviewService.\allowbreak{}decide} $\rightarrow$ \texttt{ScenarioStateMachine.\allowbreak{}apply} \\ \hline
SD-07 Gamify and play-test & \texttt{PresentationController.\allowbreak{}requestSuggestions} & \texttt{SuggestionProcessor.\allowbreak{}process}; \texttt{AssetProcessor.\allowbreak{}process} $\rightarrow$ \texttt{AssetService.\allowbreak{}checkBias}; \texttt{BuildController.\allowbreak{}buildDraftGame} \\ \hline
SD-08 Publish a scenario & \texttt{ReleaseController.\allowbreak{}publishScenario} & \texttt{BuildSnapshotter.\allowbreak{}hash} $\rightarrow$ \texttt{ReleaseService} $\rightarrow$ \texttt{StorageService.\allowbreak{}putObject}; \texttt{ReleaseScheduler.\allowbreak{}publishDueReleases} \\ \hline
SD-09 Import reference data & \texttt{ImportController.\allowbreak{}startImport} & \texttt{ImportProcessor.\allowbreak{}process} $\rightarrow$ \texttt{ReferenceSourceClient.\allowbreak{}lookup} $\rightarrow$ \texttt{NotificationService.\allowbreak{}notify} \\ \hline
SD-10 One turn of a shift & \texttt{ShiftRunController.\allowbreak{}takeTurn}, \texttt{resumeRun} & \texttt{TurnService.\allowbreak{}take} (a repeated turn identifier returns the stored result) $\rightarrow$ \texttt{NarrationProcessor.\allowbreak{}process} $\rightarrow$ \texttt{ModelGatewayService.\allowbreak{}complete} $\rightarrow$ \texttt{StreamService.\allowbreak{}publish}; for a closed segment \texttt{AssessmentService.\allowbreak{}requestGrading} $\rightarrow$ \texttt{GradingProcessor.\allowbreak{}process} $\rightarrow$ \texttt{JudgePanelService.\allowbreak{}grade}, and with too few judges \texttt{AiJobService.\allowbreak{}retry} \\ \hline
SD-11 Rate an occupation, send app feedback & \texttt{FeedbackController.\allowbreak{}sendFeedback}, \texttt{getRatingSummary}; \texttt{FeedbackInboxController.\allowbreak{}listFeedback} & \texttt{FeedbackService.\allowbreak{}send} (a rating only after the level is completed) $\rightarrow$ \texttt{FeedbackService.\allowbreak{}ratingSummary} (withheld below the threshold) $\rightarrow$ \texttt{FeedbackService.\allowbreak{}inbox} \\ \hline
SD-12 Monitor and re-run a failed job & \texttt{OperationsController.\allowbreak{}getMetrics}, \texttt{listAiCalls} & \texttt{OperationsService.\allowbreak{}metrics} (cost per run against the ceiling, latency, failed calls) $\rightarrow$ \texttt{OperationsService.\allowbreak{}aiCalls} \\ \hline
\end{longtable}
}

Description of the five AI components (Section 5.2.3).

AI-1, interviewing the domain expert. A language model asks and probes to collect what a seed needs: incidents, what a good and a poor response look like, and the facts of the situation. Each expert is interviewed separately, and the merged content marks where experts agree and where they differ. The interview is constrained by a human-\allowbreak{}written topic guide based on the critical incident technique\cite{r26} and on development practice for situational judgment tests\cite{r27}; the model probes inside each topic only. A tracker lists the seed fields still missing, the expert is shown a summary and confirms it, and every turn is stored verbatim so that each seed field links to the turn it came from. That a language model can hold an interview to a guide is supported by 766 AI-led interviews in which 94.8\% to 99.8\% of questions were open-ended, 95.7\% to 99.6\% neutral, and 0.1\% to 0.4\% hallucinated\cite{r114}.

AI-2, drafting the game seed. The model helps the content administrator turn the experts' answers into a structured seed: context, characters, situations, canon events, and the scoring frame. Drafts are written from the expert's own words, the domain expert sets every scoring key and every rubric anchor, and a seed passes automatic structure checks and then the expert's approval before any player sees it. Experts own keys and rubrics because, when GPT-4 drafted 210 examination questions, specialists required revision of 15\% and judged 0.5\% false\cite{r115}. The seed follows evidence-\allowbreak{}centered design: what is measured, which behaviour is evidence for it, and which situation elicits that behaviour\cite{r20,r116}. This is the first of four checkpoints for generated content: structure is checked by machine and content is approved by the expert.

AI-3, helping the expert test before publishing. A volunteer expert with about seven hours cannot read every free-text answer. AI-3 runs AI-4 and AI-5 on test playthroughs, from experts, testers, and simulated answers, and ranks the cases that need the expert: cases where the three judges differ by more than one level, generated turns that broke a run-time rule, and a blind random sample of cases that were not flagged. The expert decides, and can refine a rubric anchor during review, so rubrics are versioned. The random sample is required because confidence signals from a language-\allowbreak{}model grader separate right from wrong grades only weakly, with an area under the curve of 0.57 to 0.64\cite{r117}, so flagged cases alone can hide errors; if n audited cases show no error, the error rate is at most about 3/n at 95\% confidence\cite{r118}. Simulated players find relative problems only: persona-\allowbreak{}prompted models answered with near-\allowbreak{}perfect, uniform accuracy, unlike 2,359 real students\cite{r119}, so real test plays are needed as well. The run-time checks of AI-4 are to be tested on scripted player turns in four classes (on-script, off-script, adversarial, and rubric-\allowbreak{}fishing), reporting the violation rate per class; this test set is not yet built. This is the second checkpoint, and the expert-\allowbreak{}scored test plays are also the validation data for the judge. In the present phase there are no experts or real players: the team stands in for both and the advisor checks a sample, so results show that the machinery works, not agreement with experts.

AI-4, narrating during play. The game engine, not the model, controls the order of canon events. Canon content is shown verbatim from the seed; a situation may have several approved canon variants chosen at random, and the log records which variant was served. The model writes only the connecting narration, the follow-up questions that fit the player's free-text answer, and a hidden feedback note after each activity. A follow-up is decided only after the answer it follows has been graded (AI-5); a follow-up that fails the run-time checks is replaced by the expert's default follow-up, and a feedback note that fails them is redrafted, since feedback has no expert default. The whole seed is placed in the prompt, without retrieval or fine-\allowbreak{}tuning, and the rubric and scoring keys are not in the prompt. Run-time checks are rules, not another model call, so checking adds no waiting time: numbers and names must come from the seed or from the player's own text, and the engine counts follow-ups. A model from a different family reviews the logs afterwards. Holding authored structure outside free generation is the established approach in interactive narrative\cite{r120}, and when the source fits in the prompt, long-\allowbreak{}context prompting outperformed retrieval on average\cite{r121}. These are the third checkpoint (rules at run time) and the fourth (log review afterwards).

AI-5, grading each answer. The scoring model is separate from the narrating model (DC-01). There is no reference answer: answers are free text, and the judge scores each rubric criterion by how far the answer meets the behaviour described for each level. Each judge is called once per answer, follow-up answers included, with at most three criteria in a call; the call is made as soon as the answer is submitted, and the follow-up is decided only once the grade is known (FR-11, FR-12). For each criterion the judge returns a verbatim quote from the player, then the rationale, then the level, and the quote must be a substring of what the player typed. A panel of three judges from three model families grades each answer and the median is the score; cases where judges differ by more than one level go to the expert through AI-3, and if a judge is unavailable, the system switches to the other tier of the same family, so that the panel keeps three families (DC-01); if neither tier is available, the system reports an error and keeps the answers (NFR-04). One call for all criteria was not chosen because batching several rubric checks into one call trades accuracy for cost, with a large loss when the evidence is spread over a long trajectory\cite{r122}. A panel was chosen because a five-judge vote matched the best single judge, with a Pearson correlation with humans of 0.627 against 0.623\cite{r123}, so a panel reaches the best judge's level without knowing in advance which judge that is. The design does not claim that a panel is more accurate than a single judge.

Candidate models.

{\small\setstretch{1}\renewcommand{\arraystretch}{1.1}
\begin{longtable}{|>{\raggedright\arraybackslash}p{\dimexpr 0.148\textwidth-2\tabcolsep-\arrayrulewidth\relax}|>{\raggedright\arraybackslash}p{\dimexpr 0.323\textwidth-2\tabcolsep-\arrayrulewidth\relax}|>{\raggedright\arraybackslash}p{\dimexpr 0.263\textwidth-2\tabcolsep-\arrayrulewidth\relax}|>{\raggedright\arraybackslash}p{\dimexpr 0.261\textwidth-2\tabcolsep-\arrayrulewidth\relax}|}
\caption{Candidate models, with list prices, Vietnamese benchmark scores, and structured-\allowbreak{}output support.} \\
\hline
\textbf{Role} & \textbf{Candidates (USD per million tokens, input / output)} & \textbf{SEA-HELM Vietnamese} & \textbf{Schema-\allowbreak{}constrained output} \\ \hline
\endfirsthead
\hline
\textbf{Role} & \textbf{Candidates (USD per million tokens, input / output)} & \textbf{SEA-HELM Vietnamese} & \textbf{Schema-\allowbreak{}constrained output} \\ \hline
\endhead
Judge, Anthropic & claude-\allowbreak{}sonnet-\allowbreak{}5-\allowbreak{}5 (2 / 10); claude-\allowbreak{}haiku-\allowbreak{}5-\allowbreak{}5 (0.10 / 0.50) & Sonnet 5: 77.0; Haiku 4.5: 67.7 (earlier models) & Native adapter\cite{r195} \\ \hline
Judge, OpenAI & gpt-6.1-sol (2 / 10); gpt-6-luna (0.10 / 0.50) & GPT 5.6 Sol: 77.1; GPT 5.6 Luna: 71.9 (earlier models) & Strict schema \\ \hline
Judge, Google & gemini-\allowbreak{}3.\allowbreak{}8-\allowbreak{}flash (0.75 / 3.75); gemini-\allowbreak{}3.\allowbreak{}5-\allowbreak{}flash-\allowbreak{}lite (0.30 / 2.50) & gemini-\allowbreak{}3.\allowbreak{}8-\allowbreak{}flash: 78.4; Gemini 3.1 Flash Lite: 75.2 & Response schema; unsupported parameters ignored\cite{r196} \\ \hline
Narrator & qwen3.\allowbreak{}8-\allowbreak{}flash (0.15 / 0.47); deepseek-\allowbreak{}flash (0.30 / 1.20) & Qwen 3.8 Flash Next: 74.9; DeepSeek V4 Flash: 72.1 & Qwen: strict schema except in the Singapore region\cite{r200}; DeepSeek: valid JSON only, occasionally empty\cite{r201} \\ \hline
\end{longtable}
}

Threats and controls (Section 5.3.10).

{\small\setstretch{1}\renewcommand{\arraystretch}{1.1}
\begin{longtable}{|>{\raggedright\arraybackslash}p{\dimexpr 0.193\textwidth-2\tabcolsep-\arrayrulewidth\relax}|>{\raggedright\arraybackslash}p{\dimexpr 0.370\textwidth-2\tabcolsep-\arrayrulewidth\relax}|>{\raggedright\arraybackslash}p{\dimexpr 0.171\textwidth-2\tabcolsep-\arrayrulewidth\relax}|>{\raggedright\arraybackslash}p{\dimexpr 0.261\textwidth-2\tabcolsep-\arrayrulewidth\relax}|}
\caption{Threats, controls, and their status.} \\
\hline
\textbf{Threat} & \textbf{Control} & \textbf{Requirement} & \textbf{Prototype status} \\ \hline
\endfirsthead
\hline
\textbf{Threat} & \textbf{Control} & \textbf{Requirement} & \textbf{Prototype status} \\ \hline
\endhead
Stolen or guessed credentials & Passwords hashed; short-lived bearer access token, refresh token rotated on use; sign-out revokes the session; idle sessions expire & NFR-17 & Password hashing, token service, and bearer authentication exist; idle expiry not verified \\ \hline
A user acting beyond their role & Roles on every operation (\texttt{x-\allowbreak{}roles}); a second check that a staff member is assigned to the occupation concerned; guards run authentication first, then role & FR-40, FR-41 & Global authentication and role guards exist; the assignment check exists for the staff routes that are built \\ \hline
Prompt injection through player text (instructions hidden in an answer, requests to see the answer key, off-topic pressure) & The model receives the explorer's text as data inside a fixed prompt; the judge panel returns marks that must quote the answer verbatim, and a mark whose quotation is not in the answer is rejected & NFR-17, NFR-15 & Quotation check exists in the judge output parser; a scripted ``lost and attacking'' playthrough exists in the AI lab \\ \hline
Personal data leaving the country & De-\allowbreak{}identification before every model call; only seed, rubric, and de-\allowbreak{}identified text sent; allowlisted egress & DC-03, NFR-05 & De-\allowbreak{}identification exists in the narrator path; egress proxy is design only \\ \hline
Content altered after approval & Published releases immutable and versioned; every state change logged in the same transaction & NFR-16 & Lifecycle log exists; release publishing is not built (Section 6.1.3) \\ \hline
A model fails or is slow & Failover to a model of the same role; circuit breaker per provider; when none is left the run pauses and keeps its progress & NFR-04, D-92 & Gateway with failover and circuit breaker exists; a pause with ``no model'' was observed in tests \\ \hline
Runaway cost & Every call goes through the gateway into the cost ledger; ceiling per session; a budget of calls per grading & NFR-11 & Cost ledger and call budget exist; the per-session ceiling is not verified \\ \hline
Duplicate or replayed submissions & Client-\allowbreak{}generated turn identifier; one transaction per turn & NFR-03 & Idempotent answer and turn submission exist \\ \hline
Flood of requests & Web application firewall and load balancer in front; stateless pods scale out & NFR-18 & Design only; load measured on one process (Section 7.2) \\ \hline
Secrets in source or logs & Secrets from a secret manager or environment, never in the repository; errors as Problem Details without stack traces & NFR-17 & Environment validated at start-up; \texttt{.\allowbreak{}env} is not committed \\ \hline
\end{longtable}
}

Construction of the evaluation data (Section 7.3.1).

Step 1, situations. A situation is one free-text task inside a scenario, with its rubric: two or three criteria, each with a description of the behaviour at every level. Situations are taken from the provisional seeds of the three pilot roles (back-end engineering, DevOps and cloud engineering, and business analysis), which the team prepares from authoritative public sources under the advisor's guidance. They are chosen to cover the problem space on two axes the product already defines: occupation, and scenario archetype, which also fixes difficulty, since each archetype opens at a given level: execution of a clear task from L1, ambiguous requirements from L2, incidents under time pressure from L3, conflicts and reviews of others' work from L4; trade-off decisions, which open at L6, are left to Stage B. One situation per role and archetype gives 15 situations, the smallest set in which every cell of the two axes is observed. The two scenarios of the prototype supply two of them.

Step 2, responses. Two kinds of response are written for each situation. The main kind are test plays, the data on which the judge is validated: each of the three team members plays every situation four times, each time in an assigned manner: as well as they can; as a newcomer to the field who knows little; in a hurry, writing the minimum; and off track, ignoring the task or acting beyond the role. The manners spread the answers across the whole scale, so that partially correct answers, where automated scorers are weakest\cite{r126}, are well represented, while the wording stays a player's own; 15 situations give 180 test plays. The auxiliary kind are constructed responses, which test what agreement alone cannot show: a strong reference response, variants that each degrade exactly one criterion to a lower level while keeping the others, following the perturbation-\allowbreak{}checklist method\cite{r73}, and probes for what the judge must not reward or punish: a paraphrase without the rubric's wording, keyword stuffing, negation of the right action in the right words, an injected instruction to the grader, and length padding. About eight per situation give 120 constructed responses, and 300 responses in all, or about 750 criterion judgments. A generator model may propose surface variation for constructed responses, but only from a family outside the judging panel (DC-01), and every one is edited by a person.

Step 3, scores and labels. Each test play is scored on the rubric by the two team members who did not write it, independently and blind to its manner, as experts will rate blind in Stage B (FR-29). Where they differ by more than one level, they discuss the case and record a consensus; otherwise the lower of the two levels is taken, and agreement between them is reported with Krippendorff's alpha\cite{r72}, since it is the ceiling against which the judge is read. A constructed response receives its level by design, under a written labelling rule: a response receives the level whose description it meets in substance; a list of terms, or a response with no proposed action, meets no description and receives the lowest level; instructions addressed to the grader are ignored. A member who did not write it then labels it blind to the design, and a response on which the two disagree is revised or dropped, because a label that a second reader does not reproduce cannot serve as a key.

Step 4, advisor check. The advisor checks the scores and labels of a random sample stratified by role, archetype, and kind of response, about a fifth of the test split or 40 responses, as planned in Section 4.2.2. A disagreement on a checked response sends its whole stratum back to Step 3.

Step 5, freeze and split. The set is versioned and hashed. One situation per archetype, five in all, forms a development split used while prompts and rubric formats are being written; the other ten form a test split that is scored once, with frozen prompts, for the reported results. Without the split, prompts tuned on the responses that measure them would overstate the judge.

Requirement verification (Section 4.3).

The functional, non-\allowbreak{}functional, and data requirements of Section 4.3, with the business and technical verification of each.

{\small\setstretch{1}\renewcommand{\arraystretch}{1.1}
\begin{longtable}{|>{\raggedright\arraybackslash}p{\dimexpr 0.103\textwidth-2\tabcolsep-\arrayrulewidth\relax}|>{\raggedright\arraybackslash}p{\dimexpr 0.193\textwidth-2\tabcolsep-\arrayrulewidth\relax}|>{\raggedright\arraybackslash}p{\dimexpr 0.193\textwidth-2\tabcolsep-\arrayrulewidth\relax}|>{\raggedright\arraybackslash}p{\dimexpr 0.193\textwidth-2\tabcolsep-\arrayrulewidth\relax}|>{\raggedright\arraybackslash}p{\dimexpr 0.312\textwidth-2\tabcolsep-\arrayrulewidth\relax}|}
\caption{Functional requirements with business and technical verification.} \\
\hline
\textbf{ID} & \textbf{Functional requirement} & \textbf{Group} & \textbf{Business verification} & \textbf{Technical verification} \\ \hline
\endfirsthead
\hline
\textbf{ID} & \textbf{Functional requirement} & \textbf{Group} & \textbf{Business verification} & \textbf{Technical verification} \\ \hline
\endhead
FR-01 & Register, sign in with a password or Google, and recover the account & Explorer & Acceptance test: a user registers, signs in with a password and with Google, resets a forgotten password, and reaches their own progress & Authentication tests; unauthenticated requests for progress are refused; reset links expire after one use; Google sign-in verifies the ID token and finds or creates the linked account \\ \hline
FR-02 & Edit, sign out of, and delete the account & Explorer & Walkthrough of each account action & Account deletion removes personal data as defined in NFR-05 \\ \hline
FR-03 & Run an optional, retakable RIASEC orientation quiz suggesting occupations\cite{r41} & Explorer & Pilot users complete the quiz and find the suggested occupations meaningful & Quiz yields a RIASEC profile; the five suggested occupations (project-\allowbreak{}defined) come from matching it with each occupation's O*NET interest code, without AI; skipping leads to the Galaxy Map; retaking replaces the earlier result \\ \hline
FR-04 & Show the explorer's portrait for their RIASEC profile, including the social index & Explorer & Pilot users judge the description relevant and not misleading & No AI is used in the quiz; the portrait is the fixed text written in advance for the explorer's highest-\allowbreak{}scoring RIASEC type, one text per type (six types\cite{r41}); the social index is the profile's Social score; the portrait is labelled as self-\allowbreak{}reported and not validated \\ \hline
FR-05 & Browse occupations on the Galaxy Map and show each profile with its downsides and level descriptions & Explorer & Explorers find an occupation on the map and can state its main downsides and what each level involves & Every published occupation appears on the map with all profile fields and a description for every level; ratings follow FR-19 suppression \\ \hline
FR-06 & Open the entry level of every published occupation to every explorer & Explorer & Acceptance test: a new explorer opens the entry level of any published occupation & Entry level reachable with zero credit and without any adjacency check \\ \hline
FR-07 & Let explorers start an adjacent occupation at a higher level when overlap conditions are met & Explorer & Domain experts judge the offered starting levels as plausible for seeded profiles & Overlap only unlocks higher starting levels and never blocks an entry level; no offer is made where adjacency data is missing \\ \hline
FR-08 & Enrol in levels, choose scenarios, award credit, and unlock the next level & Explorer & Acceptance test: an explorer enrols, chooses scenarios, and progresses through levels as expected & Threshold boundary tests unlock at and only at the specified credit; below the threshold, tasks for the missing competencies are suggested \\ \hline
FR-09 & Offer the role, team, and context before each scenario & Explorer & Explorer walkthrough: explorers can open the context before a scenario and understand it & End-to-end test: an intro screen precedes every scenario and lets the explorer open the full context or skip it; an unavailable scenario shows a notice and returns the explorer to the Island Map \\ \hline
FR-10 & Support choice, ordering, prioritising, and free-text answers & Explorer & Explorers complete each response form without help & Known-input tests return the expected scores for each closed form \\ \hline
FR-11 & Score each free-text answer, including every follow-up answer, with anchor points backed by quotations & Explorer & Domain experts judge a sample of scored responses as fair and well-\allowbreak{}evidenced & Each answer is graded as soon as it is submitted, before any follow-up is decided, by a panel of three AI judges (project-\allowbreak{}defined) running in parallel, and their scores are combined into one; every quotation is found verbatim in the response; scores outside the rubric anchors are rejected; agreement with experts is checked under NFR-14 \\ \hline
FR-12 & Answer follow-up questions generated by AI on free-text answers & Explorer & Explorers judge the follow-up questions relevant to what they wrote & A follow-up is decided only after the previous answer is graded; follow-ups stay within the expert's follow-up goal and the activity's limit; each follow-up passes the output check before display, otherwise the expert's default follow-up is shown; each follow-up answer is graded under FR-11 \\ \hline
FR-13 & Offer hints when the explorer is stuck & Explorer & Explorers continue after a hint and still find the task meaningful & Using a hint removes the +2 option for that activity; hint use is recorded and shown in the summary; hints do not reveal grading criteria \\ \hline
FR-14 & Return result, feedback, and guidance separately & Explorer & Usability test: explorers identify the result, the feedback, and the guidance separately & All three present and distinguishable for every assessed scenario; AI-written feedback notes are built from the rubric anchors and must pass the output check (on topic, no grading criteria, correct language); a note that fails is redrafted, with no expert default; explorers can flag wrong content to the domain expert \\ \hline
FR-15 & Record evidenced expertise and soft-skill competencies, capped per skill per scenario & Explorer & Domain experts spot-check sampled sessions: recorded competencies match the observed behaviour & Records correspond to rubric criteria and are split into expertise and soft skills; random-\allowbreak{}event answers add bonus soft-skill evidence only; replay test confirms accumulation never exceeds the cap \\ \hline
FR-16 & Offer optional side quests for bonus points & Explorer & Explorers find side quests short and optional & Side quests are authored by the domain expert (FR-22); side-quest points are recorded separately from main-quest competence credit \\ \hline
FR-17 & Show the career path of each occupation level & Explorer & Explorers can describe the role held at each level of an occupation & Career path lists every level of the occupation with its role and description, as written by the domain expert \\ \hline
FR-18 & Show a private record of attempts, expertise and soft-skill points, awards, and past scenario feedback & Explorer & Explorers confirm their record is accurate and complete & Record reconstructs from session history for a seeded account; both point types shown in the top bar, summary, and profile; past feedback is the end-\allowbreak{}of-\allowbreak{}scenario feedback of each attempt; not accessible to other users \\ \hline
FR-19 & Collect occupation ratings and feedback & Explorer & Acceptance test: the rating and feedback form appears only after full completion; aggregated results are shown or withheld as expected & Form accepted only once every level is completed; aggregated result updates after each new submission; suppressed below the response-\allowbreak{}count threshold \\ \hline
FR-20 & Send app feedback and review it & Explorer, System administrator & An explorer sends feedback and a system administrator finds it in the review list & Every submission is stored with time and app version; the system administrator can mark it as handled \\ \hline
FR-21 & Manage the occupation profile: levels, level descriptions, future roles, downsides, salary, generation rules, competencies, adjacency & Domain expert & Experts approve the adaptation for Vietnamese practice & A profile can be created only for an occupation the content administrator has created; only domain experts assigned to the occupation can edit it; every rubric criterion resolves to a competency; routes without enough data are hidden, not guessed \\ \hline
FR-22 & Author scenario context, allowed task types, option scores, rubric, hints, random events, and side quests & Domain expert & Walkthrough: an author prepares a scenario for an assigned task end to end, including its hints, random events, and side quests & Edits persist with authorship recorded; access limited to domain experts assigned to that occupation; every free-text activity has a follow-up goal and a default follow-up; random events are dropped into sessions at random within the configured probability, target the scenario's reviewed skills, and never lower the main score \\ \hline
FR-23 & Generate scenario content with AI from the author's context, keeping its generation history & Domain expert & Domain-\allowbreak{}expert review of sampled generations against the required facts and the occupation's rules & Automated validator passes on every stored generation; a scenario that fails the checks cannot be submitted; each run's context, model, cost, and check result are viewable \\ \hline
FR-24 & Add reference answers with expert scores & Domain expert & Experts confirm the stored criteria reflect what they intended & Item is stored only when validation passes; failing items return the reasons \\ \hline
FR-25 & Review another expert's scenario against the review criteria & Domain expert & Walkthrough: an expert approves, rejects, and requests changes on sample scenarios using the criteria & Reviewer is never the author; every criterion is answered before a decision; unapproved scenarios are never served to explorers; a rejected scenario is retired and cannot be reopened \\ \hline
FR-26 & Play-test the draft game before publication & Domain expert & Walkthrough: the author is invited, play-tests a draft game, and reports findings & Author is notified when the draft game is built; play-test sessions are excluded from every explorer aggregate; display findings return the scenario to presentation, content findings return it to the author and to review \\ \hline
FR-27 & Revalidate a published scenario on request & Domain expert & Experts revalidate a sample of published scenarios on request & A domain expert or content administrator can send a published scenario back to the review queue; the scenario stays published and playable meanwhile; the result is recorded; outdated scenarios are flagged for replacement \\ \hline
FR-28 & Report disagreements between AI and expert grading & Domain expert & Experts confirm the disagreement list helps them refine grading criteria & Disagreement list matches recomputation over the reference set \\ \hline
FR-29 & Rate responses and scenarios blind & Domain expert & Experts confirm they cannot see AI scores or other experts' ratings while rating & Blind-\allowbreak{}rating records exist per expert; inter-\allowbreak{}expert agreement is computable \\ \hline
FR-30 & Create and manage own occupations & Content administrator & Walkthrough: a content administrator takes their occupation from draft to published to retired & State transitions follow the defined lifecycle; publication requires completed review; actions on other occupations are refused; a domain that still holds occupations cannot be deleted \\ \hline
FR-31 & Import O*NET and ESCO reference data & Content administrator & A content administrator imports a sample and adjusts it & Imported records resolve against source identifiers; unresolved records are listed; domain experts are notified to derive competencies once the import completes \\ \hline
FR-32 & Define the scenario review criteria & Content administrator & Domain experts confirm the criteria are clear enough to apply & Reviews cannot start until criteria exist for the occupation; each review records an answer per criterion \\ \hline
FR-33 & Open one active scenario per task & Content administrator & Walkthrough: a content administrator opens a scenario and the author is notified & A second active scenario for the same task is refused; retired scenarios, including rejected ones, do not count \\ \hline
FR-34 & Assign a domain expert as reviewer, with a deadline, and reassign overdue reviews & Content administrator & Walkthrough: a content administrator assigns a domain expert to review a scenario and reassigns an overdue review & Assigning the author is refused; the domain expert is notified; the deadline appears in the review queue; the reviewer is reminded before the deadline; once it passes, the content administrator is notified and can reassign the review with a new deadline \\ \hline
FR-35 & Choose the scenario presentation with AI: task types from the domain expert's allowed set, characters, and scenes & Content administrator & Walkthrough: a content administrator gamifies an approved scenario & Task types outside the domain expert's allowed set are refused; content text is read-only in this step \\ \hline
FR-36 & Maintain a bias-\allowbreak{}checked asset library; AI generates missing assets & Content administrator & Reviewers judge sampled assets, including AI-\allowbreak{}generated ones, free of stereotyping & Every asset carries its source and bias-check result; an AI-\allowbreak{}generated asset enters the library only after its bias check passes \\ \hline
FR-37 & Publish a scenario after review and play-test & Content administrator & Walkthrough: a content administrator publishes a ready scenario with an effective date and release note & Publish is disabled unless the status is ready \\ \hline
FR-38 & Delete a published scenario by retiring it & Content administrator & Walkthrough: a content administrator replaces an outdated scenario & Deleting retires the scenario and keeps its play history; the task's scenario slot is freed; retired scenarios are never served to explorers \\ \hline
FR-39 & Monitor own occupations: statistics and coverage & Content administrator & A content administrator spots a weak or outdated scenario from the dashboard & Figures match stored records; ratings follow the suppression threshold; coverage matches stored metadata \\ \hline
FR-40 & Assign and reassign one content administrator and domain experts to each occupation & System administrator & Walkthrough: a system administrator assigns staff to an occupation and moves it to another content administrator & Every occupation carries exactly one content administrator and at least one domain expert once assigned; reassigning an occupation moves its pending work to the new content administrator \\ \hline
FR-41 & Manage all user accounts, roles, and role-based access & System administrator & Walkthrough: a system administrator views the activity log of an explorer and a staff member, then disables and deletes an account & Staff sign in on the shared login screen and see only the work for their assigned occupations; wrong-role access to admin screens is refused; disabled accounts cannot sign in; deleted accounts follow NFR-05; staff-\allowbreak{}authored records are kept \\ \hline
FR-42 & Configure the generation model and the three judge models separately & System administrator & A system administrator changes a provider through the configuration screen & The three judges come from three different model families, none the same as the generation model's; configuration rejects any other combination \\ \hline
FR-43 & Monitor usage, cost, and errors, and re-run failed AI jobs & System administrator & A system administrator uses the monitoring screen to spot cost or error spikes & Metrics consistent with generated load; alerts raised when a cost threshold is crossed; failed generation and grading runs listed with cause and can be re-run \\ \hline
\end{longtable}
}

{\small\setstretch{1}\renewcommand{\arraystretch}{1.1}
\begin{longtable}{|>{\raggedright\arraybackslash}p{\dimexpr 0.114\textwidth-2\tabcolsep-\arrayrulewidth\relax}|>{\raggedright\arraybackslash}p{\dimexpr 0.182\textwidth-2\tabcolsep-\arrayrulewidth\relax}|>{\raggedright\arraybackslash}p{\dimexpr 0.212\textwidth-2\tabcolsep-\arrayrulewidth\relax}|>{\raggedright\arraybackslash}p{\dimexpr 0.210\textwidth-2\tabcolsep-\arrayrulewidth\relax}|>{\raggedright\arraybackslash}p{\dimexpr 0.277\textwidth-2\tabcolsep-\arrayrulewidth\relax}|}
\caption{Non-\allowbreak{}functional requirements with business and technical verification.} \\
\hline
\textbf{ID} & \textbf{Non-\allowbreak{}functional requirement} & \textbf{Category} & \textbf{Business verification} & \textbf{Technical verification} \\ \hline
\endfirsthead
\hline
\textbf{ID} & \textbf{Non-\allowbreak{}functional requirement} & \textbf{Category} & \textbf{Business verification} & \textbf{Technical verification} \\ \hline
\endhead
NFR-01 & Follow-up generation within 5 s (95th percentile); progress shown after 2 s & Performance & Usability session: explorers find the wait for follow-up questions acceptable and see progress feedback & Latency distribution measured over at least 200 sampled follow-ups \\ \hline
NFR-02 & Free-text assessment within 8 s (95th percentile) & Performance & Usability session: explorers find the wait for feedback acceptable & Latency distribution measured over at least 200 sampled assessments \\ \hline
NFR-03 & Progress survives interruption & Reliability & Acceptance test: an explorer closes the browser mid-\allowbreak{}scenario, returns, and continues & Forced-\allowbreak{}interruption tests resume at the recorded position \\ \hline
NFR-04 & Failover to another model when a model fails; no progress lost and user notified within 10 s when no model is available & Reliability & Explorers receive a clear message and find their progress intact & Fault injection against the provider interface: a failed model is replaced by another available model in the same role, and zero progress-\allowbreak{}record loss across injected failures \\ \hline
NFR-05 & Minimal, encrypted, deletable personal data; only de-\allowbreak{}identified answers sent to AI providers & Security and privacy & Data-\allowbreak{}protection review against applicable regulation; walkthrough of an account deletion & Data inventory review; encryption check; account deletion removes records; AI requests contain no name, e-mail, or account ID; providers are configured with no training on submitted data and no retention \\ \hline
NFR-06 & Suppress aggregates below a minimum population & Security and privacy & Privacy review confirms small-\allowbreak{}population figures are withheld and the reason is stated & Aggregates below threshold are withheld \\ \hline
NFR-07 & Disclose AI generation and AI assessment & Transparency & Pilot users can state that scenarios and assessments are machine-\allowbreak{}produced & Disclosure present at the point of use in every relevant screen \\ \hline
NFR-08 & Vietnamese-\allowbreak{}language interface and content & Usability & Native speakers review the interface and content for completeness and accuracy & Every screen and all content render in Vietnamese with no missing text \\ \hline
NFR-09 & At least 80\% of readers state the task objective correctly & Usability & Comprehension check with a minimum of 20 representative readers & Validator confirms every scenario states an explicit objective in its context \\ \hline
NFR-10 & New occupations need no code change & Maintainability & A content administrator publishes a new occupation through the authoring interface alone & No code change or redeployment is recorded for the new occupation \\ \hline
NFR-11 & Session cost stays within the configured ceiling & Cost & System administrator reviews cost against the agreed budget & Cost telemetry aggregated per session and compared against the configured ceiling \\ \hline
NFR-12 & Model provider replaceable without content changes & Portability & Domain experts confirm scenario quality is unchanged after a provider switch & Provider substitution test with unchanged content; regression suite passes \\ \hline
NFR-13 & Works in a mobile browser on a slow connection & Compatibility & Explorers complete a scenario on their own phone & Test on reference devices and a throttled network profile agreed for testing \\ \hline
NFR-14 & AI grading validated against expert judgment & Fairness & Domain experts accept the agreement report for their occupation & Quadratic weighted kappa per difficulty stratum recomputed whenever a judge model changes \\ \hline
NFR-15 & Published scenarios match approved content; AI content during play stays within expert-set limits & Accuracy & Reviewers confirm sampled published scenarios and sampled follow-ups and feedback stay within the approved content & Automated comparison between approved and published content shows no difference; every follow-up shown passed the output check or is the expert's default follow-up; every feedback note shown passed the output check \\ \hline
NFR-16 & Scenario status changes logged; published content immutable & Auditability & A content administrator traces who approved and published a scenario & Lifecycle log complete for every scenario; write attempts on published content are refused \\ \hline
NFR-17 & Passwords hashed, all traffic encrypted, idle sessions expire, AI prompts protected against injection & Security and privacy & Security review of sign-in, session handling, and AI prompts & Stored passwords are salted hashes; plain-HTTP requests are redirected or refused; sessions end after the configured idle time; explorer answers reach the AI as separated data whose instructions are ignored; an injection test set in the regression suite never changes a score \\ \hline
NFR-18 & Explorer service available within the agreed operating target & Availability & A system administrator reviews monthly availability against the agreed target & Uptime monitoring records availability; planned maintenance is announced in advance \\ \hline
NFR-19 & Explorer screens meet WCAG 2.1 level AA & Accessibility & Users with assistive technology complete a scenario & Automated and manual accessibility audit against WCAG 2.1 AA \\ \hline
\end{longtable}
}

{\small\setstretch{1}\renewcommand{\arraystretch}{1.1}
\begin{longtable}{|>{\raggedright\arraybackslash}p{\dimexpr 0.103\textwidth-2\tabcolsep-\arrayrulewidth\relax}|>{\raggedright\arraybackslash}p{\dimexpr 0.210\textwidth-2\tabcolsep-\arrayrulewidth\relax}|>{\raggedright\arraybackslash}p{\dimexpr 0.204\textwidth-2\tabcolsep-\arrayrulewidth\relax}|>{\raggedright\arraybackslash}p{\dimexpr 0.208\textwidth-2\tabcolsep-\arrayrulewidth\relax}|>{\raggedright\arraybackslash}p{\dimexpr 0.269\textwidth-2\tabcolsep-\arrayrulewidth\relax}|}
\caption{Data requirements with source, business verification, and technical verification.} \\
\hline
\textbf{ID} & \textbf{Data requirement} & \textbf{Source} & \textbf{Business verification} & \textbf{Technical verification} \\ \hline
\endfirsthead
\hline
\textbf{ID} & \textbf{Data requirement} & \textbf{Source} & \textbf{Business verification} & \textbf{Technical verification} \\ \hline
\endhead
DR-01 & Occupational reference data (occupations, tasks, skills) & O*NET and ESCO, imported by the content administrator and adapted for Vietnamese practice by domain experts & Domain experts review the adaptation for Vietnamese practice & Imported records resolve against source identifiers \\ \hline
DR-02 & Competency taxonomy with observable indicators & Derived from DR-01 under expert adaptation & Domain experts review definitions and indicators & Every rubric criterion resolves to a taxonomy entry \\ \hline
DR-03 & Occupational adjacency with overlap measures & Computed from DR-01/DR-02 & Domain experts review computed adjacency for plausibility & Overlap values are reproducible from DR-01/DR-02 \\ \hline
DR-04 & Scenario content: context, activities, option scores, rubric, hints, random events, side quests & Author domain expert with AI assistance & Reviewer sign-off recorded before presentation & Context schema validation; automatic checks pass before submission \\ \hline
DR-05 & Generated scenarios with provenance & System & A content administrator traces any served scenario back to its context and model & Every served scenario carries resolvable provenance \\ \hline
DR-06 & User progress: enrolled levels, competence credit capped per scenario, side-quest points, play history, end-\allowbreak{}of-\allowbreak{}scenario feedback, awards & System & Explorers confirm their progress and awards are shown correctly & Progress reconstructs from stored records; replaying a scenario never adds credit beyond the cap \\ \hline
DR-07 & Assessment records, including each judge's score and the combined grade per answer, hidden feedback notes, hint use, and explorer flags & System and AI judges, with flags from explorers & Domain experts audit a sample of assessment records and review flagged content & Records complete for every assessed response, follow-up answers included; every feedback note stores its output-\allowbreak{}check result; every flag reaches a domain expert of that occupation \\ \hline
DR-08 & Five-\allowbreak{}attribute occupation ratings (project-\allowbreak{}defined) & Explorers & Acceptance test of the rating flow after full completion & Ratings exist only for fully completed occupations \\ \hline
DR-09 & Expert reference set stratified by difficulty & Domain experts & Inter-\allowbreak{}expert agreement recorded for the reference set & Set covers the defined difficulty strata; every item passed validation \\ \hline
DR-10 & Operational telemetry & System & A system administrator answers cost and error questions from the telemetry & Telemetry present and queryable \\ \hline
DR-11 & Play-test session flag & System & Staff confirm dashboards are unchanged after play-\allowbreak{}testing & Play-test sessions are absent from explorer-\allowbreak{}facing aggregates \\ \hline
DR-12 & Quiz answers, RIASEC profile with social index, and the fixed portrait texts (one per RIASEC type) & Explorers; portrait texts written in advance by the project team & Pilot users judge the portrait texts relevant and not misleading & Profile recalculates from stored answers; every RIASEC type has a portrait text; retaking the quiz replaces the profile and portrait \\ \hline
DR-13 & Scenario lifecycle record, including author and reviewer assignments with review deadlines, review decisions, comments, and revalidation results & System, from staff actions & A content administrator reconstructs a scenario's history & Every transition in SM-01 is logged; every assignment, reassignment, review, and revalidation result is stored with actor and time; at most one active scenario per task \\ \hline
DR-14 & Presentation data and asset library & Content administrator with AI & Reviewers check sampled assets for stereotyping & Every presented asset resolves to a library entry with a bias-check result \\ \hline
DR-15 & Occupation catalogue (domains and occupations) and occupation profiles (level descriptions, future roles, downsides) & Content administrator (catalogue) and domain experts (profiles) & A content administrator reviews the catalogue; domain experts review their occupation's profile for completeness & Every published occupation has all required fields and a description for every level; every profile field records the domain expert who wrote it \\ \hline
DR-16 & User accounts, linked Google identities, roles, assignments, and activity log & System administrator; activity log from the system & A system administrator confirms each occupation's staff and reviews an account's activity & Exactly one content administrator and at least one domain expert per occupation; a Google identity links to one account only; every account action is logged \\ \hline
DR-17 & Platform settings & System administrator & A system administrator reviews settings against policy & Every change is versioned with actor and time \\ \hline
DR-18 & App feedback records & Explorers & A system administrator reviews recent feedback & Each record holds text, time, app version, and handling status \\ \hline
DR-19 & Scenario review criteria & Content administrator & Domain experts confirm the criteria match how they review & Criteria are versioned; each review links to the version it used \\ \hline
\end{longtable}
}

Architecture details (Section 5.3.1).

The in-flight figure follows from Little's law: 85 calls per second, each lasting three to eight seconds. Two sizing assumptions remain to be measured by the load tests of Chapter 7: if one single-vCPU API replica serves about 300 light requests per second, Stage 4 needs four replicas, six with headroom across two zones; if one worker holds 50 calls in flight, it needs six to fourteen workers.

The architecture scales for seven reasons. Processing tiers are stateless, so capacity grows by adding replicas. Fast and slow work are separated: a request stores and acknowledges in about 0.2 s, and slow work waits in a queue, so a class of forty starting together lengthens the queue rather than overloading the API. The binding constraint, the model tier, is isolated and limited on its own. Published content is immutable and therefore almost entirely served from cache. Reads for statistics are separated from writes by a replica. Load is shed by priority rather than by failure: under pressure, live grading and narration go first, feedback notes wait a little longer for the summary screen, and authoring work waits for quiet hours. And no transaction spans services, so growth never requires distributed coordination. Canon content is approved once and served verbatim from releases, so run-time model calls are limited to narration, follow-ups, feedback notes, and grading; the cost per playthrough is estimated in Section 5.2.3.

Reliability and operations. Availability of 99.5\% a month is proposed for Stages 1 and 2 and 99.9\% once two zones run (NFR-18). The write-ahead log is archived continuously and backups taken nightly, both copied to a second region in Vietnam, for a proposed recovery point of five minutes and recovery time of one hour. Deployments roll replica by replica behind health checks, with backward-\allowbreak{}compatible migrations so old and new versions run side by side. Monitoring tracks queue length, 95th-\allowbreak{}percentile latency per provider, cost per session, and error rates, alerting at the thresholds of NFR-01, NFR-02, and NFR-11. A web application firewall and TLS guard the entry, secrets are held in a secret manager, and the egress proxy admits only allowlisted model endpoints.

Modules and tables (Section 5.3.4).

{\small\setstretch{1}\renewcommand{\arraystretch}{1.1}
\begin{longtable}{|>{\raggedright\arraybackslash}p{\dimexpr 0.249\textwidth-2\tabcolsep-\arrayrulewidth\relax}|>{\raggedright\arraybackslash}p{\dimexpr 0.148\textwidth-2\tabcolsep-\arrayrulewidth\relax}|>{\raggedright\arraybackslash}p{\dimexpr 0.483\textwidth-2\tabcolsep-\arrayrulewidth\relax}|>{\raggedright\arraybackslash}p{\dimexpr 0.114\textwidth-2\tabcolsep-\arrayrulewidth\relax}|}
\caption{NestJS modules, the processes that run them, the tables they write, and their routes.} \\
\hline
\textbf{Module} & \textbf{Processes} & \textbf{Writes tables of} & \textbf{Routes} \\ \hline
\endfirsthead
\hline
\textbf{Module} & \textbf{Processes} & \textbf{Writes tables of} & \textbf{Routes} \\ \hline
\endhead
AuthModule & API & Group A except assignment, through the one account store & 7 \\ \hline
UserModule & API & Assignment of staff to occupations (\texttt{ASSIGNED\_\allowbreak{}TO}) & 10 \\ \hline
CatalogueModule & API & Groups B and G (occupations, soft skills, tiers, NPC cards, outfits), loaded by a data-\allowbreak{}loading command & 6 \\ \hline
ImportModule & Worker, API & O*NET and ESCO import, \texttt{KEYWORD} & 4 \\ \hline
ScenarioModule & Worker, API & Groups D and E, expert interviews, generation runs & 16 \\ \hline
LifecycleModule & Shared & \texttt{LIFECYCLE\_\allowbreak{}LOG}; the SM-01 state machine & 0 \\ \hline
ReviewModule & API & Review criteria, reviews, revalidation & 10 \\ \hline
PublishingModule & Worker, API & Assets, presentation, builds, play-test findings, releases & 12 \\ \hline
CareerModule & API & Groups H and I & 9 \\ \hline
LifeModule & API & Group K and side quests & 6 \\ \hline
PlayModule & Worker, API & Runs and run events (group J) & 11 \\ \hline
AssessmentModule & Worker & Assessments, judge results, feedback notes & 0 \\ \hline
QualityModule & Worker, API & Content flags, blind ratings, regrading & 6 \\ \hline
AiModule & API & Prompts, models, jobs, model calls (group L) & 2 \\ \hline
FeedbackModule & API & Occupation ratings, app feedback & 4 \\ \hline
MonitoringModule & API & None; reads usage, cost, and errors & 3 \\ \hline
CoreModule & Global, shared & Activity log, platform settings & 0 \\ \hline
NotificationModule & Global, worker, API & Notifications; the event stream & 3 \\ \hline
GameCoreModule & Global, shared & None; a snapshot of reference tables built once at start-up for \textit{game-\allowbreak{}core}, which never reads the database & 0 \\ \hline
\end{longtable}
}

Entries awaiting a decision (Section 5.3.5).

Entries still waiting for a decision. Some entities are drawn from research that is still marked as a draft or as to be decided: \texttt{DEPARTMENT}, the outfit entities, \texttt{CHECK\_\allowbreak{}QUESTION}, \texttt{CONDUCT\_\allowbreak{}FLAG}, \texttt{SKILL\_\allowbreak{}TIER} and its tier names, \texttt{ACHIEVEMENT}, \texttt{DISTRESS\_\allowbreak{}TERM}, \texttt{CAREER\_\allowbreak{}HABIT}, \texttt{LIFE\_\allowbreak{}MILESTONE}, and \texttt{SCENARIO.\allowbreak{}outcomes}. Prices of shop items, the amount of health recovered, and the cost of replaying a scenario have no figure yet, so tests check the relations between them, not their values. The turn identifier with which the explorer application makes a repeated submission safe (SD-10) has no attribute of \texttt{RUN\_\allowbreak{}EVENT} yet; the prototype keeps it in the player-\allowbreak{}input event. Two readings of the diagram also need confirming: the \texttt{(min,max)} labels count the partners of the neighbouring entity in 13 relationships and the entity itself in the rest, and some lookup tables are drawn as having at least one row on the other side. They are items A28 and A29 of the open issues. The prototype's database does not follow this diagram yet, so migrating from it is part of the implementation (Chapter 6).

Judge strength and reliability (Section 5.2.3).

Strength of the judge relative to the generator. The judge does not need to be stronger than the generator, because it scores text written by a person, not text the generator wrote. It does need enough capacity: on German clinical-\allowbreak{}reasoning transcripts, a small model agreed poorly with human raters (intraclass correlation below 0.20) and a larger one moderately (0.41 to 0.57), against 0.45 between the two human raters\cite{r124}. Whether a larger tier, or a reasoning model, is worth its cost for this task is not settled by published evidence, so each judge family is tested at two tiers and the cheaper tier is kept unless it is clearly worse.

Reliability of a language-\allowbreak{}model judge. Published evidence supports a conditional answer, not reliability by default, and none of it concerns Vietnamese or role-play. A synthesis of 65 studies of essay scoring found agreement between language models and human raters to vary widely and to depend on context\cite{r101}. Untrained general models fall well short of dedicated scorers: on the ASAP essays, ChatGPT and Llama reached an average quadratic weighted kappa of 0.313 and 0.297, against 0.817 for a supervised system\cite{r125}. Agreement comes close to human level only under conditions: GPT-4 given a few scored examples reached a kappa of 0.70 on short answers in school subjects, against 0.75 between humans\cite{r102}, and approached dedicated scoring systems on CEFR essays once calibration examples were supplied\cite{r103}. Three weaknesses recur. Scorers handle fully correct and fully incorrect answers well and degrade on partially correct ones, where human agreement stays stable\cite{r126}. Without a reference answer, judges agree with experts mainly on questions they can answer themselves\cite{r127}, and a model's accuracy as a judge closely follows its accuracy at solving the same problems\cite{r60}; JobQuest has no reference answers, so the judge must know IT work, and the level descriptions of the rubric must carry the domain knowledge. And consistency is not validity: GPT-4's repeated ratings agreed with each other at an intraclass correlation of 0.94 to 0.99\cite{r105}, which says nothing about agreement with people. The design answers each weakness: a strong model and an explicit rubric from the start, scored examples later from expert rulings, a panel whose disagreements send partially correct cases to the expert, and level descriptions written by the expert. The project then measures reliability on its own task instead of assuming it (RP-2), against the criteria of Williamson, Xi and Breyer\cite{r71}: quadratic weighted kappa of at least 0.70, a standardised mean difference of at most 0.15, and agreement at most 0.10 below that between experts. The data, raters, measures, and thresholds are fixed in Section 7.3.

Items awaiting game design. Several items cannot be settled before the game design is fixed: the scoring scale, where $-$1 / 0 / +1 / +2 is planned and not final and the current sample scenarios describe three levels only; the competency scoring table for a playthrough and the specific criteria given to the judge; the list of required facts in a seed, needed to measure fact preservation (NFR-15); and how evidence from the play process, such as hints, time, and follow-ups, counts and how scores combine into a competency. The AI design does not depend on them: the judge scores any criterion that comes with level descriptions, on any ordinal scale (Section 5.3.6), and agreement with experts can be computed on any such scale.

Components (Section 5.3.1).

Explorer application. The occupation map, levels, the scenario player for the four response forms of FR-10, hints, follow-ups, results, the personal record, ratings, and feedback (FR-03 to FR-20). It runs \textit{game-core} for instant provisional feedback on closed formats, queues unsent turns locally, and receives narration, follow-ups, and assessment results through a server-sent event stream. Prototype: screens, the 3D map, and the four forms exist; the turn queue, event stream, and 2D map do not.

Staff application. The thirty-\allowbreak{}three staff use cases: the AI-1 interview through which a domain expert gives the knowledge a seed needs, review, test playthroughs, reliability, and blind rating for domain experts; seed drafting with AI-2, catalogue, presentation, publication, and statistics for content administrators; accounts, model configuration, monitoring, and feedback for system administrators. It shares the sign-in screen and origin with the explorer application (FR-41) but is a separate download. Prototype: absent.

Game core. A library with no framework, network, or storage dependency: the reducer folding run events into state, scoring of closed formats against expert keys, credit with the per-\allowbreak{}scenario cap, unlocking and adjacency offers, content schemas, and the rule checks on model output of Section 5.2.3 --- numbers and names in narration found in the seed or the player's text, the follow-up count, each judge's quotation found verbatim in the answer, levels within the scale. It runs unchanged in browser and server. Prototype: present; its free-text grader is a synchronous keyword matcher inside the reducer, which must become an asynchronous assessment entering the run as an event.

Identity and access. Accounts, sign-in and recovery, sessions, occupation-\allowbreak{}scoped roles, disabling, deletion, and the activity log (FR-01, FR-02, FR-40, FR-41); every module asks it whether an actor may act on an occupation. Prototype: registration, email and Google sign-in, rotating sessions, and deletion exist; reset, idle expiry, scoped roles, disabling, and the log do not, and the three prototype roles must become the four actors of Section 4.2.1.

Exploration. Everything outside the assessment itself: orientation, catalogue and profiles, enrolment, runs and their events, credit, unlocking, adjacency offers, side quests, ratings, records, and feedback (FR-03 to FR-09, FR-15 to FR-20). It validates turn events with \textit{game-core} and stores them; each free-text answer is handed to Assessment, which stores it with an assessment request in the same transaction. Prototype: present, but runs are submitted as one batch at the end.

Assessment. Turns answers into results, feedback, and guidance (FR-10 to FR-14). Closed formats are scored synchronously by \textit{game-\allowbreak{}core}; each free-text answer, follow-up answers included, is graded by the judge panel as soon as it is submitted (AI-5), and only then is a follow-up decided (FR-12). The grade and each activity's checked feedback note are written back as events, and at the end of the scenario the notes are combined into the feedback shown on the summary screen in the three separate parts of FR-14. It also routes explorers' flags of wrong content to the occupation's domain experts. Prototype: closed formats scored by replaying the run; free text matched by keyword.

Content lifecycle. Occupations, knowledge, scenarios and drafts, generation history, review criteria, assignments and decisions, presentation, assets, and publication (FR-21 to FR-38). It is the only writer of content, implements the state machine SM-01, writes each transition with its log entry in one transaction, and publishes immutable, versioned releases. It uses AI orchestration for generation, validation, presentation, and missing assets; the scheduler for deadlines and annual revalidation; and the importer for O*NET and ESCO, whose interest codes also drive orientation matching (FR-03). A domain expert or content administrator can send a published scenario back for revalidation while it stays playable (FR-27). Prototype: versioned immutable documents, written only by a seed script.

AI orchestration. Runs every model use of the five AI components of Section 5.2.3 as a job: the expert interview (AI-1) and seed drafting (AI-2) at authoring time, with presentation design and missing assets (FR-35, FR-36); test playthroughs (AI-3); narration, follow-ups, and feedback notes during play (AI-4, FR-12, FR-14); and grading of each answer by a panel of three judges (AI-5, FR-11). The Get-\allowbreak{}to-\allowbreak{}know-\allowbreak{}me quiz uses no AI (FR-04). Per job it loads the pinned content version, removes identifying details from explorer text, checks the session budget, calls the model port, applies \textit{game-\allowbreak{}core}'s rule checks, and records model, prompt version, cost, and latency. There is no scripted substitute for a model's output: a job that cannot complete is retried or waits, and a session pauses until it can. It rechecks at job time that no judge shares the narrator's family or another judge's. Prototype: absent; queue and model-\allowbreak{}client libraries are declared but unused.

Model gateway. One adapter per provider and a registry recording each model's role --- text generation, assessment, image generation (FR-36) --- with its provider and family; configuration is refused when the narrator and a judge, or two judges, share either (FR-42, DC-01). It applies per-\allowbreak{}provider concurrency limits, timeouts, retries, and a circuit breaker. Adapters are written by the project over the OpenAI-\allowbreak{}compatible format, with native adapters where that format falls short (Section 5.2.3). Prototype: absent.

Analytics. Everything shown only in aggregate: ratings and statistics above the population threshold, play-tests excluded (FR-19, FR-39, NFR-06, DR-11); coverage; agreement by difficulty stratum and the disagreement list (FR-28, NFR-14); and blind ratings (FR-29). Occupation ratings are re-\allowbreak{}aggregated in the transaction that stores each rating (FR-19); heavier statistics refresh periodically, later from a replica. It never returns an individual's record. Prototype: personal summaries only.

Operations. Telemetry of usage, cost, latency, and errors (DR-10); failed jobs and re-runs (FR-43); versioned settings, including the cost ceiling (DR-17); and feedback review (FR-20). Prototype: absent.

Notification, scheduler, and import. Notification sends email and in-app notices: reset links, review assignments, reminders, play-test invitations. The scheduler runs recurring jobs: review reminders and overdue notices (FR-34), annual revalidation (BR-15), retention, refreshing statistics, and the reconciliation that re-queues any narration or scoring request left without a result. The importer loads O*NET and ESCO and lists unresolved records (FR-31). Prototype: absent.

Data stores. The operational database is the system of record: relational tables for accounts, assignments, runs, events, evidence, reviews, ratings, and logs, and documents for content. The queue store holds jobs, rate limits, budget counters, and the publish--subscribe channel feeding event streams --- nothing that cannot be rebuilt from the database. The object store holds assets and exports. Prototype: the database, with 19 tables.

Interface principles and screens (Section 5.3.8).

{\small\setstretch{1}\renewcommand{\arraystretch}{1.1}
\begin{longtable}{|>{\raggedright\arraybackslash}p{\dimexpr 0.220\textwidth-2\tabcolsep-\arrayrulewidth\relax}|>{\raggedright\arraybackslash}p{\dimexpr 0.126\textwidth-2\tabcolsep-\arrayrulewidth\relax}|>{\raggedright\arraybackslash}p{\dimexpr 0.649\textwidth-2\tabcolsep-\arrayrulewidth\relax}|}
\caption{Design principles of the interface and their sources.} \\
\hline
\textbf{Principle} & \textbf{Source} & \textbf{Consequence for the interface} \\ \hline
\endfirsthead
\hline
\textbf{Principle} & \textbf{Source} & \textbf{Consequence for the interface} \\ \hline
\endhead
All text in Vietnamese, with no language selection & NFR-08, D-96 & Interface strings are written directly in Vietnamese; there is no translation layer to leave a screen half-\allowbreak{}translated \\ \hline
Usable on a phone over a slow connection & NFR-13, DC-05 & Layouts are designed from the narrow width up; the 3D map and its libraries load only when the 3D view opens; a budget of 200 KB for the first interactive screen (Section 5.2.1) \\ \hline
Accessible without a pointer or a canvas & NFR-19 & Every planet of the Galaxy Map is also a focusable, named element in an SVG view that can be chosen at any time and is used automatically when WebGL is missing, when reduced motion is requested, or on low-memory devices \\ \hline
Result, feedback, and guidance are three separate things & FR-14 & The summary screen shows them in separate blocks, never as one score with a comment \\ \hline
Waiting is never silent & NFR-01, NFR-02, NFR-04 & The answer screen shows ``saved, assessing'' at once, then progress after two seconds; when no model is available it says so and keeps the run \\ \hline
Machine-\allowbreak{}produced content is labelled & NFR-07 & A disclosure appears at the point of use wherever a narration, a follow-up, or an assessment is shown \\ \hline
Small populations are not shown & NFR-06 & Aggregated ratings show the reason they are withheld instead of a figure \\ \hline
Content shown is content the expert approved & NFR-15, SM-01 & The play screens show only text of a published release or of a play-test draft; there is no free-text field through which the explorer can change a scenario \\ \hline
\end{longtable}
}

{\small\setstretch{1}\renewcommand{\arraystretch}{1.1}
\begin{longtable}{|>{\raggedright\arraybackslash}p{\dimexpr 0.350\textwidth-2\tabcolsep-\arrayrulewidth\relax}|>{\raggedright\arraybackslash}p{\dimexpr 0.441\textwidth-2\tabcolsep-\arrayrulewidth\relax}|>{\raggedright\arraybackslash}p{\dimexpr 0.204\textwidth-2\tabcolsep-\arrayrulewidth\relax}|}
\caption{Screens of the explorer application.} \\
\hline
\textbf{Screen} & \textbf{Route} & \textbf{Requirements} \\ \hline
\endfirsthead
\hline
\textbf{Screen} & \textbf{Route} & \textbf{Requirements} \\ \hline
\endhead
Welcome, sign in, register, forgot and reset password & \texttt{/\allowbreak{}welcome}, \texttt{/\allowbreak{}login}, \texttt{/\allowbreak{}register}, \texttt{/\allowbreak{}forgot-\allowbreak{}password}, \texttt{/\allowbreak{}reset-\allowbreak{}password} & FR-01 \\ \hline
Galaxy Map (home) & \texttt{/\allowbreak{}galaxy}, \texttt{/\allowbreak{}home} & FR-05 \\ \hline
Get-\allowbreak{}to-\allowbreak{}know-\allowbreak{}me quiz and its result & \texttt{/\allowbreak{}quiz}, \texttt{/\allowbreak{}quiz/\allowbreak{}result} & FR-03, FR-04 \\ \hline
Occupation profile and level selection & \texttt{/\allowbreak{}occupations/\allowbreak{}:code}, \texttt{/\allowbreak{}occupations/\allowbreak{}:code/\allowbreak{}levels} & FR-05, FR-06, FR-08, FR-17 \\ \hline
Island map and onboarding of a scenario & \texttt{/\allowbreak{}islands/\allowbreak{}:enrolmentId}, \texttt{/\allowbreak{}islands/\allowbreak{}:enrolmentId/\allowbreak{}onboard} & FR-08, FR-09, FR-16 \\ \hline
Play & \texttt{/\allowbreak{}play/\allowbreak{}:runId} & FR-10 to FR-13, FR-15 \\ \hline
Reflection and summary & \texttt{/\allowbreak{}play/\allowbreak{}:runId/\allowbreak{}reflect}, \texttt{/\allowbreak{}play/\allowbreak{}:runId/\allowbreak{}summary} & FR-14, FR-15 \\ \hline
Profile & \texttt{/\allowbreak{}profile} & FR-18 \\ \hline
Rate occupation & \texttt{/\allowbreak{}occupations/\allowbreak{}:code/\allowbreak{}rate} & FR-19 \\ \hline
App feedback & \texttt{/\allowbreak{}feedback} & FR-20 \\ \hline
Account settings & \texttt{/\allowbreak{}settings} & FR-02 \\ \hline
\end{longtable}
}

{\small\setstretch{1}\renewcommand{\arraystretch}{1.1}
\begin{longtable}{|>{\raggedright\arraybackslash}p{\dimexpr 0.384\textwidth-2\tabcolsep-\arrayrulewidth\relax}|>{\raggedright\arraybackslash}p{\dimexpr 0.426\textwidth-2\tabcolsep-\arrayrulewidth\relax}|>{\raggedright\arraybackslash}p{\dimexpr 0.185\textwidth-2\tabcolsep-\arrayrulewidth\relax}|}
\caption{Screens of the staff application.} \\
\hline
\textbf{Screen} & \textbf{Route (under \texttt{/\allowbreak{}staff})} & \textbf{Requirements} \\ \hline
\endfirsthead
\hline
\textbf{Screen} & \textbf{Route (under \texttt{/\allowbreak{}staff})} & \textbf{Requirements} \\ \hline
\endhead
Catalogue of occupations; reference-\allowbreak{}data import & \texttt{catalogue}, \texttt{catalogue/\allowbreak{}imports/\allowbreak{}:importId?} & FR-30, FR-31 \\ \hline
Scenarios of an occupation; review criteria & \texttt{occupations/\allowbreak{}:occupationId/\allowbreak{}scenarios}, \texttt{.\allowbreak{}.\allowbreak{}.\allowbreak{}/\allowbreak{}review-\allowbreak{}criteria} & FR-32, FR-33, FR-34 \\ \hline
Author a scenario: context, generation, reference answers & \texttt{scenarios/\allowbreak{}:scenarioId/\allowbreak{}author} & FR-22, FR-23, FR-24 \\ \hline
Review a scenario & \texttt{scenarios/\allowbreak{}:scenarioId/\allowbreak{}review} & FR-25 \\ \hline
Gamify: choose presentation, characters, scenes & \texttt{scenarios/\allowbreak{}:scenarioId/\allowbreak{}gamify} & FR-35 \\ \hline
Asset library & \texttt{assets} & FR-36 \\ \hline
Publish & \texttt{scenarios/\allowbreak{}:scenarioId/\allowbreak{}publish} & FR-37 \\ \hline
System prompts of an occupation & \texttt{occupations/\allowbreak{}:occupationId/\allowbreak{}prompts} & FR-23 \\ \hline
Monitor an occupation & \texttt{occupations/\allowbreak{}:occupationId/\allowbreak{}monitor} & FR-39 \\ \hline
Settings & \texttt{settings} & --- \\ \hline
\end{longtable}
}

Third parties (Section 5.3.10).

{\small\setstretch{1}\renewcommand{\arraystretch}{1.1}
\begin{longtable}{|>{\raggedright\arraybackslash}p{\dimexpr 0.161\textwidth-2\tabcolsep-\arrayrulewidth\relax}|>{\raggedright\arraybackslash}p{\dimexpr 0.224\textwidth-2\tabcolsep-\arrayrulewidth\relax}|>{\raggedright\arraybackslash}p{\dimexpr 0.302\textwidth-2\tabcolsep-\arrayrulewidth\relax}|>{\raggedright\arraybackslash}p{\dimexpr 0.159\textwidth-2\tabcolsep-\arrayrulewidth\relax}|>{\raggedright\arraybackslash}p{\dimexpr 0.148\textwidth-2\tabcolsep-\arrayrulewidth\relax}|}
\caption{Third parties and the data they receive.} \\
\hline
\textbf{Third party} & \textbf{Purpose} & \textbf{What it receives} & \textbf{Where} & \textbf{Personal data} \\ \hline
\endfirsthead
\hline
\textbf{Third party} & \textbf{Purpose} & \textbf{What it receives} & \textbf{Where} & \textbf{Personal data} \\ \hline
\endhead
Content delivery network & Serve both applications and published releases & Static bundles and releases & Edge locations, including abroad & None \\ \hline
Google Identity & Optional sign-in & An identity token, only when the user chooses Google & Abroad & Yes, by the user's choice \\ \hline
E-mail provider & Password resets and notices & Address and message & To be chosen with a host in Vietnam & Yes \\ \hline
Model providers & Narration, scoring, feedback, authoring help & Seed content and de-\allowbreak{}identified answer text, under an opaque run identifier & Abroad & Not by design; residual risk below \\ \hline
O*NET, ESCO & Reference data, imported & Nothing; data flows in only & --- & None \\ \hline
Second region of the cloud provider & Backups and recovery & Encrypted copies of the database and logs & Vietnam & Yes, encrypted \\ \hline
\end{longtable}
}

Frameworks of career guidance (Section 2.4).

{\small\setstretch{1}\renewcommand{\arraystretch}{1.1}
\begin{longtable}{|>{\raggedright\arraybackslash}p{\dimexpr 0.184\textwidth-2\tabcolsep-\arrayrulewidth\relax}|>{\raggedright\arraybackslash}p{\dimexpr 0.092\textwidth-2\tabcolsep-\arrayrulewidth\relax}|>{\raggedright\arraybackslash}p{\dimexpr 0.193\textwidth-2\tabcolsep-\arrayrulewidth\relax}|>{\raggedright\arraybackslash}p{\dimexpr 0.322\textwidth-2\tabcolsep-\arrayrulewidth\relax}|>{\raggedright\arraybackslash}p{\dimexpr 0.203\textwidth-2\tabcolsep-\arrayrulewidth\relax}|}
\caption{Frameworks of career decision difficulty compared with the taxonomy of Gati, Krausz and Osipow.} \\
\hline
\textbf{Framework} & \textbf{Year} & \textbf{What it describes} & \textbf{Structure} & \textbf{Missing occupational information as a separate component} \\ \hline
\endfirsthead
\hline
\textbf{Framework} & \textbf{Year} & \textbf{What it describes} & \textbf{Structure} & \textbf{Missing occupational information as a separate component} \\ \hline
\endhead
Taxonomy of career decision-\allowbreak{}making difficulties (CDDQ)\cite{r1} & 1996 & Difficulties arising before and during the decision & Ten categories in three clusters: lack of readiness, lack of information, inconsistent information & Yes: within the lack-\allowbreak{}of-\allowbreak{}information cluster, alongside information about the self, the process, and ways of obtaining information \\ \hline
Emotional and personality-\allowbreak{}related career decision-\allowbreak{}making difficulties (EPCD)\cite{r106} & 2008 & Emotional and personality-\allowbreak{}related difficulties & Three clusters: pessimistic views, anxiety, self-\allowbreak{}concept and identity & No \\ \hline
Career Indecision Profile\cite{r107} & 2013 & Sources of indecision, including personality & Four factors: neuroticism and negative affectivity, choice and commitment anxiety, lack of readiness, interpersonal conflicts & No separate component \\ \hline
Career Decision Self-\allowbreak{}Efficacy scale\cite{r108} & 1983 & Confidence in carrying out decision tasks & Five subscales: self-\allowbreak{}appraisal, occupational information, goal selection, planning, problem solving & Partly: confidence in gathering information, not the information that is lacking \\ \hline
Career Thoughts Inventory\cite{r109} & 1996 & Dysfunctional career thoughts & Three scales: decision-\allowbreak{}making confusion, commitment anxiety, external conflict & No \\ \hline
Career Decision Scale\cite{r110} & 1976 & Certainty and indecision & A certainty scale and an indecision scale & No separate component \\ \hline
\end{longtable}
}
